% ==============================================================================
% PART II: SYNTACTIC CONSTITUENTS & ANTI-TOP-HEAVY GRAMMAR
% File: chapters/part2_syntax.tex
% ==============================================================================
\part{Syntactic Architecture \& Anti-Top-Heavy Grammar / 句法骨架与防头重脚轻语法}\label{part:syntax}

\chapter{Five Foundational Sentence Patterns in Scholarly Writing / 学术写作五大核心基本句型}\label{ch:patterns}
一切复杂的学术长难句，本质上均由五大基本句型（Five Canonical Clause Patterns）经过扩展、修饰后置与从属复合演化而来。掌握句子核心主干（Subject, Predicate Verb, Object, Complement）的定位与成分配比，是杜绝中式流水句与残缺句的第一道防线。

\section{Constituent Foundations \& Clause Architecture / 句子骨架与核心成分解析}
英语是典型的“主干驱动型（Head-Driven）”语言。一个合格的学术句子必须具备清晰的主谓核心：
\begin{itemize}[leftmargin=1.5em]
  \item \textbf{主语 (Subject, S)}: 承载核心学术实体，力求紧凑凝练（建议 $\le 5$ 词）。
  \item \textbf{谓语动词 (Predicate Verb, V)}: 决定句型骨架与后置补足成分配价（Valency）。
  \item \textbf{宾语 (Object, O)}: 动作的直接或间接承受者。
  \item \textbf{补足语 (Complement, C)}: 对主语或宾语状态、性质进行定性定量的评测补充。
  \item \textbf{状语 (Adverbial, A)}: 交代时间、空间、因果、条件、程度等上下文边界。
\end{itemize}

\section{The Five Canonical Clause Patterns in Research Literature / 五大基本句型学术深度应用}
\subsection{Pattern I: SV (Subject + Intransitive Verb) --- 自然过程与客观演变}
\begin{ConceptBox}{Structural Breakdown \& Academic Utility}
SV 句型由不及物动词引导，常用于客观描述系统自发变化、自然衰减或数学收敛过程。核心谓语动词后可后置介词短语状语。
\end{ConceptBox}
\begin{ExampleBox}{Publication-Grade Exemplar}
\textit{The training loss converges rapidly during the initial five epochs.}
\end{ExampleBox}
\vspace{4pt}
\subsection{Pattern II: SVO (Subject + Transitive Verb + Object) --- 主动干预与实验操作}
\begin{ConceptBox}{Structural Breakdown \& Academic Utility}
SVO 是学术写作中最坚实、最常见的主动干预句型。主语代表模型、算法或作者团队，谓语动词直接施加动作于客体。
\end{ConceptBox}
\begin{ExampleBox}{Publication-Grade Exemplar}
\textit{The transformer architecture captures long-range temporal dependencies across sensor channels.}
\end{ExampleBox}
\vspace{4pt}
\subsection{Pattern III: SVC (Subject + Linking Verb + Complement) --- 属性归属与状态界定}
\begin{ConceptBox}{Structural Breakdown \& Academic Utility}
SVC 用于刻画客体的内在性质。系动词（remain, prove, appear, be of）连接主语与表语，表语赋予主语学术价值评估。
\end{ConceptBox}
\begin{ExampleBox}{Publication-Grade Exemplar}
\textit{The estimated parameter remains stable across varied noise levels.}
\end{ExampleBox}
\vspace{4pt}
\subsection{Pattern IV: SVOO (Subject + Dative Verb + Indirect Object + Direct Object) --- 资源分配与功能赋予}
\begin{ConceptBox}{Structural Breakdown \& Academic Utility}
双宾语句型体现资源供给或机制赋予，通常可转换为 SVO + Preposition 结构（provides X for Y / gives X to Y）。
\end{ConceptBox}
\begin{ExampleBox}{Publication-Grade Exemplar}
\textit{The platform provides researchers with an integrated evaluation environment.}
\end{ExampleBox}
\vspace{4pt}
\subsection{Pattern V: SVOC (Subject + Complex Transitive + Object + Complement) --- 批判性裁定与效能定性}
\begin{ConceptBox}{Structural Breakdown \& Academic Utility}
SVOC 是学术评价中极具表现力的高阶句型。宾语补足语（suboptimal）直接定性动作导致的结果状态，论证力量极其强劲。
\end{ConceptBox}
\begin{ExampleBox}{Publication-Grade Exemplar}
\textit{Ablation trials render the initial baseline suboptimal under heavy computation load.}
\end{ExampleBox}
\vspace{4pt}
\chapter{Principle of End-Weight and Information Distribution / 尾重原则与信息焦点分布}\label{ch:endweight}
初学者在学术写作中最致命的毛病就是“头重脚轻（Top-Heavy Subject）”：将一个长达十几词的动名词短语或主语从句堆在句首，谓语动词极度延后，迫使读者长时间挂起短时认知记忆，产生极其沉重的阅读负荷。本章系统讲解认知语言学中的尾重原则（End-Weight Principle）以及三项外科重构手术。

\section{Cognitive Processing Load \& The Anti-Top-Heavy Imperative / 认知负荷与防头重脚轻原理}
英语母语读者的视觉聚焦模型遵循“短主语启动 $\to$ 强谓语连接 $\to$ 复杂信息尾部汇聚”的认知节奏。若句首主语超过 8 个词，读者大脑在定位到核心动作前已产生认知疲劳。遵循尾重原则即是将长短语、定语从句及复杂条件后移至谓语动词之后。

\begin{ConceptBox}{The Golden Rule of End-Weight (尾重黄金法则)}
在学术正式书面语中，轻量主干居前，重量信息居后：\par
\begin{center}
\small\texttt{[Compact Subject ($\le$5 words)] + [Active Verb] + [Heavy Object / Clause / Prepositional Modifiers]}
\end{center}
\end{ConceptBox}

\section{Three Reconstructive Surgeries for Top-Heavy Subjects / 头重脚轻病句的三大外科重构手术}
\subsection{Surgery 1: Extraposition with Anticipatory 'It' (形式主语 it 替代法)}
\begin{ExampleBox}{Defective Top-Heavy vs. Balanced Reconstruction}
\textbf{\color{TargetHighlight}[Top-Heavy Defect]}:\par\textit{That modern deep reinforcement learning policies suffer from severe distribution shifts when deployed in real robots is widely recognized.}\par\vspace{4pt}
\textbf{\color{AcadNavy}[End-Weight Reconstruction]}:\par\textit{It is widely recognized that modern deep reinforcement learning policies suffer from severe distribution shifts when deployed in real robots.}\par\vspace{4pt}
\textbf{Surgical Rationale}: 将沉重的主语从句 That ... 移至句尾，句首使用轻巧的形式主语 it 占位，主谓核心 it is widely recognized 在前 4 词内迅速建立。
\end{ExampleBox}
\vspace{4pt}
\subsection{Surgery 2: Active Reporting Subject Extraction (强施事主干提取法)}
\begin{ExampleBox}{Defective Top-Heavy vs. Balanced Reconstruction}
\textbf{\color{TargetHighlight}[Top-Heavy Defect]}:\par\textit{The observation that the gradient magnitude increases non-linearly with batch size was made by the team during training.}\par\vspace{4pt}
\textbf{\color{AcadNavy}[End-Weight Reconstruction]}:\par\textit{The team observed that gradient magnitude increases non-linearly with batch size during training.}\par\vspace{4pt}
\textbf{Surgical Rationale}: 将隐藏在介词短语或被动语态中的真实施事者（The team）提取为全句主动语态主语，动词由 was made 升华为直接宾语从句引导词 observed。
\end{ExampleBox}
\vspace{4pt}
\subsection{Surgery 3: Inversion \& Heavy-NP Shift (完全倒装与长成分后移法)}
\begin{ExampleBox}{Defective Top-Heavy vs. Balanced Reconstruction}
\textbf{\color{TargetHighlight}[Top-Heavy Defect]}:\par\textit{Several fundamental limitations regarding out-of-distribution generalization under non-stationary sensor dynamics lie at the heart of this dispute.}\par\vspace{4pt}
\textbf{\color{AcadNavy}[End-Weight Reconstruction]}:\par\textit{At the heart of this dispute lie several fundamental limitations regarding out-of-distribution generalization under non-stationary sensor dynamics.}\par\vspace{4pt}
\textbf{Surgical Rationale}: 将句末简短的地点介词短语 At the heart of this dispute 置于句首充当逻辑主位，谓语动词 lie 紧随其后，超长并列主语后置平稳落地。
\end{ExampleBox}
\vspace{4pt}
\section{The Given-to-New Contract at the Clause Level / 句内“已知---新信息”尾端汇聚法则}
尾重原则与篇章的“已知---新信息契约（Given-to-New Contract）”高度统一：已知信息作为简短主位引领句首，新信息包含大量技术细节与修饰，自然居于句尾承重。

\chapter{Academic Grammar Mastery: The 36 Golden Rules (G001--G036) / 学术语法全景图谱：36条黄金法则}\label{ch:grammar36}
学术语法绝非机械死板的条条框框，而是确保科学交流准确、无歧义且具备严密学术严谨度的工程规约。本章完整收录标准学术体系的 36 条黄金语法法则（G001--G036），分为六大进阶教学集群，逐条提供原则阐释、正反例对比与警示。

\section{Cluster 1: Clause Architecture \& Basic Mechanics (G001--G006) / 句法架构与核心骨干}
\Needspace{6\baselineskip}
\subsection{G001: 五大基本句型（SV, SVO, SVC, SVOO, SVOC）与动词配价边界}\label{sec:rule:G001}
\begin{ConceptBox}{Core Grammar Principle \& Theoretical Framework}
\textbf{Category}: \texttt{C1\_basic\_clause\_patterns}\par\vspace{3pt}
\textbf{Principle}: 英语句子的核心骨架由谓语动词的及物性与配价决定：不及物动词构成 SV（Evidence accumulates / A problem arose），系动词构成 SVC（The label is meaningless），单及物动词构成 SVO（The team abandoned the approach），双及物动词构成 SVOO，复合宾语动词构成 SVOC（The pronoun makes the final sentence ambiguous）。严禁把不及物动词当及物动词用，或把及物动词误加介词。\par\vspace{4pt}
\textbf{Operational Rules}:
\begin{itemize}[leftmargin=1.5em, itemsep=1pt]
  \item 不及物动词（arise, emerge, occur, differ, fluctuate, accumulate, oscillate）后不直接接名词宾语，需用状语或介词短语扩展。
  \item 单及物动词（address, discuss, access, affect, investigate, assess, warrant）直接接宾语，禁止按中文习惯误加介词（禁止写 *discuss about, *address to, *access to 作动词）。
  \item SVOC 结构中宾语补足语用形容词或分词说明宾语状态（keep the room temperature constant, make the task easier）。
\end{itemize}
\end{ConceptBox}
\begin{ExampleBox}{Authentic Examples \& Contrasts}
\textit{SV: A problem arose when the data formats changed.}\par\vspace{2pt}
\textit{SVO: This paper addresses unequal access to educational resources.}\par\vspace{2pt}
\end{ExampleBox}
\begin{TipBox}{Diagnostic Warning \& Common Errors}
\begin{itemize}[leftmargin=1.5em, itemsep=1pt]
  \item 误将不及物动词用于被动或直接加宾语：*A problem was arisen / *differ two groups。
  \item 在及物动词后添加冗余介词：*discuss about the findings / *address to the limitation。
  \item SVOC 中误用副词作宾补：*keep the room temperature constantly（应为 constant）。
\end{itemize}
\end{TipBox}
\vspace{6pt}
\subsection{G002: 可数名词、不可数名词及多义词在不同义项下的可数性切换}\label{sec:rule:G002}
\begin{ConceptBox}{Core Grammar Principle \& Theoretical Framework}
\textbf{Category}: \texttt{C4\_A2\_nouns\_countability\_across\_senses}\par\vspace{3pt}
\textbf{Principle}: 不仅要区分绝对不可数名词（evidence, information, research, literature, equipment, apparatus, documentation, feedback, knowledge）与可数名词（study, finding, discrepancy, anomaly, covariate），更要掌握“同一单词在不同词义下的可数性切换”（Rule A.2）。\par\vspace{4pt}
\textbf{Operational Rules}:
\begin{itemize}[leftmargin=1.5em, itemsep=1pt]
  \item 不可数学术名词不加 a/an 或复数 -s（严禁写 *an evidence, *researches, *equipments, *feedbacks）；表单项用 a piece/item of evidence, a research study, a piece of equipment。
  \item 一词多义下的可数性切换：discipline 表“纪律”不可数，表“学科”可数（two academic disciplines）；work 表“研究工作”不可数，workings 表“计算步骤”；experience 表“经验”不可数，表“经历”可数（experiences in childhood）；paper 表“纸”不可数，表“论文”可数。
  \item 量词精确匹配：修饰不可数用 a great deal of, a large amount of, much, less and less；修饰可数复数用 a large number of, a handful of, many, fewer and fewer。
\end{itemize}
\end{ConceptBox}
\begin{ExampleBox}{Authentic Examples \& Contrasts}
\textit{The transcription required a great deal of effort.}\par\vspace{2pt}
\textit{A large number of records still require checking.}\par\vspace{2pt}
\end{ExampleBox}
\begin{TipBox}{Diagnostic Warning \& Common Errors}
\begin{itemize}[leftmargin=1.5em, itemsep=1pt]
  \item 量词与可数性错配：*a great deal of participants / *an amount of respondents / *fewer time。
  \item 不可数名词误加复数：*many evidences / *new equipments / *useful feedbacks。
\end{itemize}
\end{TipBox}
\vspace{6pt}
\subsection{G003: 主谓一致、数量短语一致与邻近名词干扰排除}\label{sec:rule:G003}
\begin{ConceptBox}{Core Grammar Principle \& Theoretical Framework}
\textbf{Category}: \texttt{C2\_subject\_verb\_agreement}\par\vspace{3pt}
\textbf{Principle}: 复杂主语的主谓一致取决于核心名词（Head Noun）而非后置介词短语中的邻近名词。尤其需区分 a number of 与 the number of，以及 along with / accompanied by / as well as 连接的附加成分。\par\vspace{4pt}
\textbf{Operational Rules}:
\begin{itemize}[leftmargin=1.5em, itemsep=1pt]
  \item the number of + 复数名词作主语，谓语用单数；a (large/small) number of + 复数名词作主语，谓语用复数。
  \item 主语后接 along with, together with, accompanied by, as well as 时，谓语动词与前面的核心主语保持一致。
  \item the bulk of / a proportion of / a percentage of + 名词作主语时，谓语单复数取决于 of 后的名词是不可数/单数还是复数可数名词。
  \item means, series 单复数同形，谓语由限定词决定（a series of observations was vs two series were）。
\end{itemize}
\end{ConceptBox}
\begin{ExampleBox}{Authentic Examples \& Contrasts}
\textit{The number of completed interviews was recorded separately.}\par\vspace{2pt}
\textit{A small number of respondents left this item blank.}\par\vspace{2pt}
\end{ExampleBox}
\begin{TipBox}{Diagnostic Warning \& Common Errors}
\begin{itemize}[leftmargin=1.5em, itemsep=1pt]
  \item 被 of 短语中的复数名词误导：*The analysis of interview responses show（应为 shows）。
  \item 把 as well as 误当成 and 改变主语单复数。
\end{itemize}
\end{TipBox}
\vspace{6pt}
\subsection{G004: 冠词系统（a/an, the, 零冠词）与学术特指/泛指边界}\label{sec:rule:G004}
\begin{ConceptBox}{Core Grammar Principle \& Theoretical Framework}
\textbf{Category}: \texttt{C4\_articles\_and\_specificity}\par\vspace{3pt}
\textbf{Principle}: 冠词决定名词短语的指称范围：单数可数泛指用 a/an；指代本研究特定样本/模型（the sample, the model）或带 of 后置限定的属性（the accuracy of the measurements, the nature of the task）用 the；复数可数或不可数表普遍概念用零冠词。\par\vspace{4pt}
\textbf{Operational Rules}:
\begin{itemize}[leftmargin=1.5em, itemsep=1pt]
  \item 单数可数名词在句中不可裸奔（bare singular），必须带冠词或限定词（a limitation, the protocol, each trajectory）。
  \item 注意冠词改变词义：nature（零冠词）指“大自然”，而 the nature of the problem 指“问题的基本性质”；in charge of（主动负责） vs in the charge of（由他人看管负责）。
  \item 图表编号直接作专名指称时不加 the（see Table 1, Figure 2 shows），但普通名词需加冠词（in the table below）。
\end{itemize}
\end{ConceptBox}
\begin{ExampleBox}{Authentic Examples \& Contrasts}
\textit{Equity is an abstract concept that requires an operational definition.}\par\vspace{2pt}
\textit{The nature of the task determines which measures are appropriate.}\par\vspace{2pt}
\end{ExampleBox}
\begin{TipBox}{Diagnostic Warning \& Common Errors}
\begin{itemize}[leftmargin=1.5em, itemsep=1pt]
  \item 单数可数名词漏掉冠词：*We conducted survey / *Trained model achieved lower error。
  \item 图表编号前误加冠词：*See the Table 2（应为 See Table 2）。
\end{itemize}
\end{TipBox}
\vspace{6pt}
\subsection{G005: 词形变化、词缀构词与词性语法槽位匹配（名、动、形、副）}\label{sec:rule:G005}
\begin{ConceptBox}{Core Grammar Principle \& Theoretical Framework}
\textbf{Category}: \texttt{A4\_word\_forms\_and\_morphology}\par\vspace{3pt}
\textbf{Principle}: 掌握学术词族与常见词缀（-able, -cy, -ify, -ise/-ize, -ism/-ist, -less, anti-, co-, il-/im-/in-/ir-, multi-, over-, pre-, post-）的转换，确保名词充当主宾表、动词充当谓语、形容词修饰名词或作表语、副词修饰动词/形容词/全句。\par\vspace{4pt}
\textbf{Operational Rules}:
\begin{itemize}[leftmargin=1.5em, itemsep=1pt]
  \item 严禁词性混填：accurate (adj.) / accuracy (n.) / accurately (adv.)；analyse (v.) / analysis (n., 复数 analyses) / analytical (adj.)；vary (v.) / variable (n./adj.) / variation (n.)。
  \item 警惕词缀语义陷阱：predictable 表示“可预测的”而非“已被预测的”；successive 表示“连续相继的”而非“成功的（successful）”；priceless/meaningless 语义受语境限定；bi-monthly 有歧义，学术写作应改写为 every two months 或 twice a month。
  \item 英美拼写（analyse/analyze, behaviour/behavior, parametrisation/parameterization）在同一篇文稿中必须统一。
\end{itemize}
\end{ConceptBox}
\begin{ExampleBox}{Authentic Examples \& Contrasts}
\textit{Accurate records are essential for checking the analysis.}\par\vspace{2pt}
\textit{Attendance increased in three successive years.}\par\vspace{2pt}
\end{ExampleBox}
\begin{TipBox}{Diagnostic Warning \& Common Errors}
\begin{itemize}[leftmargin=1.5em, itemsep=1pt]
  \item 形容词与名词混用：*assess the accurate of measurements（应为 accuracy）。
  \item 混淆同根异义派生词：把 successive（连续的）误作 successful（成功的），把stationary（平稳的）误拼为 stationery（文具）。
\end{itemize}
\end{TipBox}
\vspace{6pt}
\subsection{G006: 核心介词规约搭配与名/动/形+介词固定结构}\label{sec:rule:G006}
\begin{ConceptBox}{Core Grammar Principle \& Theoretical Framework}
\textbf{Category}: \texttt{B\_collocations\_and\_prepositions}\par\vspace{3pt}
\textbf{Principle}: 学术词汇的介词搭配高度固定，且同一词族的名词、动词和形容词常要求不同介词（如 access a database [v.] vs access to services [n.] vs accessible to users [adj.]；emphasise the need [v.] vs place emphasis on [n.]）。\par\vspace{4pt}
\textbf{Operational Rules}:
\begin{itemize}[leftmargin=1.5em, itemsep=1pt]
  \item 动词+介词固定搭配：account for, attribute A to B, base A on B, confine A to B, consist of, contribute to, correlate with, decompose A into B, depend on, differ from, disentangle A from B, draw on, focus on, incorporate A into B, rely on, stem from。
  \item 名词/形容词+介词精确方向：change in (某指标变化) vs change of (更换)；increase/decrease in (指标) vs of (数值)；attributable to, akin to, amenable to, analogous to, compatible with, consistent with, invariant under, proportional to, superior to, synonymous with。
  \item 当 to 为介词时（an approach to, contribute to, be devoted to, prior to, with regard to, amenable to, receptive to, subject to），后接动词必须用动名词（-ing）！
\end{itemize}
\end{ConceptBox}
\begin{ExampleBox}{Authentic Examples \& Contrasts}
\textit{The team developed an approach to handling missing observations.}\par\vspace{2pt}
\textit{The simplified model is amenable to analytical treatment.}\par\vspace{2pt}
\end{ExampleBox}
\begin{TipBox}{Diagnostic Warning \& Common Errors}
\begin{itemize}[leftmargin=1.5em, itemsep=1pt]
  \item 把介词 to 误当成不定式符号：*an approach to handle / *contribute to improve / *be devoted to solve。
  \item 介词错配：*superior than（应为 superior to）、*synonymous to（应为 synonymous with）、*attributable for（应为 attributable to）。
\end{itemize}
\end{TipBox}
\vspace{6pt}
\section{Cluster 2: Tense, Aspect, Voice \& Reference (G007--G012) / 时态语态与指代基石}
\Needspace{6\baselineskip}
\subsection{G007: 时态、体与学术论文各章节的方法论时间分工}\label{sec:rule:G007}
\begin{ConceptBox}{Core Grammar Principle \& Theoretical Framework}
\textbf{Category}: \texttt{C3\_tense\_aspect\_time}\par\vspace{3pt}
\textbf{Principle}: 学术论文的时态传递方法论状态：一般过去时描述本研究特定实验步骤与过去采样结果；一般现在时描述数学定义、图表展示内容、模型架构属性及本文篇章导引；现在完成时总结某领域迄今的研究演进或前文已完成论证。\par\vspace{4pt}
\textbf{Operational Rules}:
\begin{itemize}[leftmargin=1.5em, itemsep=1pt]
  \item 描述本研究特定数据集采集、预处理步骤与单次实验观测用一般过去时（The team abandoned the approach; Two analysts analysed the data independently）。
  \item 描述数学性质、算法机制、图表内容及普遍事实用一般现在时（The affine function f(x) = ax + b includes an intercept b; Table 1 shows...）。
  \item 描述近十年研究趋势或前文已展现的内容用现在完成时（The field has developed in recent years; As we have seen, the two estimates use different denominators）。
\end{itemize}
\end{ConceptBox}
\begin{ExampleBox}{Authentic Examples \& Contrasts}
\textit{The team accessed the database before the policy changed.}\par\vspace{2pt}
\textit{A weakly stationary process has a constant finite mean and autocovariance that depends only on the lag.}\par\vspace{2pt}
\end{ExampleBox}
\begin{TipBox}{Diagnostic Warning \& Common Errors}
\begin{itemize}[leftmargin=1.5em, itemsep=1pt]
  \item 把尚未开展的 Future Work 写成已完成的一般过去时。
  \item 在明确过去时间状语（during the pilot）下误用现在完成时。
\end{itemize}
\end{TipBox}
\vspace{6pt}
\subsection{G008: 语序、否定辖域、部分否定与句首否定倒装}\label{sec:rule:G008}
\begin{ConceptBox}{Core Grammar Principle \& Theoretical Framework}
\textbf{Category}: \texttt{C2\_word\_order\_and\_negation}\par\vspace{3pt}
\textbf{Principle}: 精确控制否定与限定副词（only, solely, merely, by no means, under no circumstances, rarely, seldom）的辖域。当否定或唯一下界短语置于句首时，主句必须采用部分倒装。\par\vspace{4pt}
\textbf{Operational Rules}:
\begin{itemize}[leftmargin=1.5em, itemsep=1pt]
  \item 区分全部否定（none of, in no way, under no circumstances, by no means）与部分否定（not all, not necessarily, not always）。
  \item 当 by no means, in no way, under no circumstances, in no circumstances, on no account, rarely, seldom, only when... 置于句首时，主句必须部分倒装（助动词/be/情态动词提至主语前）。
  \item 间接疑问从句（assess to what extent the results depend on...）一律保持陈述语序，不倒装。
\end{itemize}
\end{ConceptBox}
\begin{ExampleBox}{Authentic Examples \& Contrasts}
\textit{By no means is the result certain. (对比句中：A large sample is by no means sufficient to eliminate systematic bias.)}\par\vspace{2pt}
\textit{Under no circumstances should test labels be used for model tuning.}\par\vspace{2pt}
\end{ExampleBox}
\begin{TipBox}{Diagnostic Warning \& Common Errors}
\begin{itemize}[leftmargin=1.5em, itemsep=1pt]
  \item 句首否定强调漏掉倒装：*Under no circumstances test labels should be used。
  \item 间接疑问句误用倒装：*examine how does the system behave（应为 how the system behaves）。
\end{itemize}
\end{TipBox}
\vspace{6pt}
\subsection{G009: 代词指代一致性、比较替代词（that/those of）与消除指代悬空}\label{sec:rule:G009}
\begin{ConceptBox}{Core Grammar Principle \& Theoretical Framework}
\textbf{Category}: \texttt{C5\_pronouns\_and\_reference}\par\vspace{3pt}
\textbf{Principle}: 代词必须具有唯一、明确的先行词。学术写作中严禁出现可同时指向多个名词的歧义代词（ambiguous pronoun）或缺乏总结名词的裸奔 this。\par\vspace{4pt}
\textbf{Operational Rules}:
\begin{itemize}[leftmargin=1.5em, itemsep=1pt]
  \item 当句中有两个同数名词时，不用裸 it/they/this，应改用 the former / the latter 或 This + 总结性范畴名词（this discrepancy, this constraint, these observations）。
  \item 比较结构中必须用 that of 替代前文单数/不可数名词，用 those of 替代前文复数可数名词，确保比较对象同范畴（The cost is low relative to that of the earlier procedure）。
\end{itemize}
\end{ConceptBox}
\begin{ExampleBox}{Authentic Examples \& Contrasts}
\textit{The pronoun makes the final sentence ambiguous.}\par\vspace{2pt}
\textit{The cost is low relative to that of the earlier procedure.}\par\vspace{2pt}
\end{ExampleBox}
\begin{TipBox}{Diagnostic Warning \& Common Errors}
\begin{itemize}[leftmargin=1.5em, itemsep=1pt]
  \item 比较结构漏掉 that/those of，导致拿“误差”直接跟“基线模型”相比：*The error of our model is lower than the baseline（应为 lower than that of the baseline）。
\end{itemize}
\end{TipBox}
\vspace{6pt}
\subsection{G010: 不定式与动名词的补足语分工及非谓语动词搭配}\label{sec:rule:G010}
\begin{ConceptBox}{Core Grammar Principle \& Theoretical Framework}
\textbf{Category}: \texttt{C7\_nonfinite\_infinitives\_gerunds}\par\vspace{3pt}
\textbf{Principle}: 动词后接不定式（to do）还是动名词（doing）取决于动词配价规约；同时注意 warrant, facilitate, advocate, involve, consider, anticipate 等核心学术动词的补足语限制。\par\vspace{4pt}
\textbf{Operational Rules}:
\begin{itemize}[leftmargin=1.5em, itemsep=1pt]
  \item 接动名词（doing）或名词、不接不定式的典型动词：consider, involve, advocate, anticipate, recommend, suggest, warrant（warrant further investigation / warrant revising），facilitate（facilitate comparison，严禁写 *facilitate sb to do）。
  \item 接不定式（to do）的典型动词：aim, attempt, tend, seek, intend, resolve, enable sb to do, allow sb to do, force sb to do, oblige sb to do。
\end{itemize}
\end{ConceptBox}
\begin{ExampleBox}{Authentic Examples \& Contrasts}
\textit{The report advocates providing clearer information to applicants.}\par\vspace{2pt}
\textit{The procedure involves checking every record twice.}\par\vspace{2pt}
\end{ExampleBox}
\begin{TipBox}{Diagnostic Warning \& Common Errors}
\begin{itemize}[leftmargin=1.5em, itemsep=1pt]
  \item 误用不定式：*warrant to conclude / *facilitate researchers to compare / *advocate to provide / *involve to check。
\end{itemize}
\end{TipBox}
\vspace{6pt}
\subsection{G011: 主动语态与被动语态的基本结构、时态/情态变形与及物性限制}\label{sec:rule:G011}
\begin{ConceptBox}{Core Grammar Principle \& Theoretical Framework}
\textbf{Category}: \texttt{D\_active\_and\_passive\_voice}\par\vspace{3pt}
\textbf{Principle}: 被动语态由 be + 过去分词构成，可结合各种时态与情态动词。转换时必须保持原句时态、体、情态、否定和逻辑关系，且严禁将不及物动词或状态短语强行变为被动。\par\vspace{4pt}
\textbf{Operational Rules}:
\begin{itemize}[leftmargin=1.5em, itemsep=1pt]
  \item 只有及物动词或带受事宾语的短语动词（analyse, allocate, carry out, account for, base on, expose to, confine to）才能转换为被动语态，且短语介词不可遗漏。
  \item 不及物动词及状态短语绝不可用于被动语态：consist of, occur, arise, emerge, happen, belong to, differ, appear, exist, take place, stem from, result from, coincide with（严禁写 *is consisted of, *was occurred, *was belonged to）。
\end{itemize}
\end{ConceptBox}
\begin{ExampleBox}{Authentic Examples \& Contrasts}
\textit{Active only: The questionnaire consists of three sections. (绝不可写 *is consisted of)}\par\vspace{2pt}
\textit{Passive with modal: An anomalous observation should be investigated before it is removed.}\par\vspace{2pt}
\end{ExampleBox}
\begin{TipBox}{Diagnostic Warning \& Common Errors}
\begin{itemize}[leftmargin=1.5em, itemsep=1pt]
  \item 将不及物动词被动化：*The index is consisted of three parts / *A problem was arisen / *The pattern was emerged。
\end{itemize}
\end{TipBox}
\vspace{6pt}
\subsection{G012: 数值、比例、分位数、矩阵量纲与增减幅度的精确语法}\label{sec:rule:G012}
\begin{ConceptBox}{Core Grammar Principle \& Theoretical Framework}
\textbf{Category}: \texttt{A1\_C4\_numerical\_and\_quantitative\_grammar}\par\vspace{3pt}
\textbf{Principle}: 定量描述必须严格区分“差值（by）”、“终值（to）”、“比例（proportion/percentage）”、“比值（ratio）”、“百分点（percentage points）”以及“均值（mean）、中位数（median）、众数（mode）、中程数（midrange）”。\par\vspace{4pt}
\textbf{Operational Rules}:
\begin{itemize}[leftmargin=1.5em, itemsep=1pt]
  \item 动词 + by + 差值（decreased by 5 units） vs 动词 + to + 终值（decreased to 15）；名词 + of + 变化量（growth of 10\%） vs 名词 + in + 变化对象（growth in demand）。
  \item 从 20\% 升至 25\% 是 an increase of 5 percentage points（绝对百分点变化），而相对增幅为 25\%（relative percentage change）。
  \item double 表示增至两倍，halve 表示减半（不可写 *halve by half）；midrange 是 (min+max)/2，median 是排序后居中的分位数（0.5 quantile），二者绝不可混同。
\end{itemize}
\end{ConceptBox}
\begin{ExampleBox}{Authentic Examples \& Contrasts}
\textit{The count increased by ten, from twenty to thirty.}\par\vspace{2pt}
\textit{A rise from 20\% to 25\% is five percentage points.}\par\vspace{2pt}
\end{ExampleBox}
\begin{TipBox}{Diagnostic Warning \& Common Errors}
\begin{itemize}[leftmargin=1.5em, itemsep=1pt]
  \item 介词 by 与 to 混用导致差值与终值颠倒。
  \item 把 percentage point（百分点）与 per cent（百分比）混同。
\end{itemize}
\end{TipBox}
\vspace{6pt}
\section{Cluster 3: Lexical Precision, Hedging \& Modifiers (G013--G018) / 词汇精度、谨慎陈述与修饰从句}
\Needspace{6\baselineskip}
\subsection{G013: 近义词与多义词深度辨析：语义范围、褒贬色彩、技术定义与语域}\label{sec:rule:G013}
\begin{ConceptBox}{Core Grammar Principle \& Theoretical Framework}
\textbf{Category}: \texttt{A1\_A5\_A6\_A7\_A8\_lexical\_and\_polysemy\_precision}\par\vspace{3pt}
\textbf{Principle}: 区分近义词与一词多义在语义范围、判断强度、褒贬色彩及学科专有定义上的差别。禁止把仅有部分重合的词当作完全同义词替换，禁止为堆砌大词而改变事实或判断强度。\par\vspace{4pt}
\textbf{Operational Rules}:
\begin{itemize}[leftmargin=1.5em, itemsep=1pt]
  \item 测量与统计核心辨析：accuracy（准确度，接近真值） vs precision（精密度，重复测量一致）；reliable（稳定一致） vs valid（有效测量目标构念或推理成立）；efficacy（受控试验效力） vs effectiveness（真实场景效果）；equal（均等） vs equitable（按需求公平分配）；seasonality（固定日历周期） vs cyclicality（非固定周期的波动）；isotonic（保序的/等渗的） vs isotropic（各向同性的）。
  \item 褒贬与语域辨析：simple（简明清晰的） vs simplistic（过度简化的，贬义）；exploit structure（中性：利用结构） vs exploit workers（贬义：剥削）；manipulate a variable（中性：实验操纵自变量） vs falsify data（贬义：伪造数据）。
\end{itemize}
\end{ConceptBox}
\begin{ExampleBox}{Authentic Examples \& Contrasts}
\textit{An equitable allocation of resources may differ from an equal allocation.}\par\vspace{2pt}
\textit{Business cyclicality may involve expansions and contractions of varying duration, whereas seasonality follows a fixed calendar period.}\par\vspace{2pt}
\end{ExampleBox}
\begin{TipBox}{Diagnostic Warning \& Common Errors}
\begin{itemize}[leftmargin=1.5em, itemsep=1pt]
  \item 把 suggest/indicate 机械替换为 prove/demonstrate 导致过度推断。
  \item 混淆形近学术词：adapt/adopt, adverse/averse, discrete/discreet, temporal/temporary, stationary/stationery, isotonic/isotropic。
\end{itemize}
\end{TipBox}
\vspace{6pt}
\subsection{G014: 情态动词、确定性梯度与科学谨慎表达（Epistemic Hedging）}\label{sec:rule:G014}
\begin{ConceptBox}{Core Grammar Principle \& Theoretical Framework}
\textbf{Category}: \texttt{C6\_E\_modals\_and\_hedging}\par\vspace{3pt}
\textbf{Principle}: 科学写作要求结论的语气强度与实验/数据证据的完备度严格对齐。通过情态动词（may, might, can, could）、认知副词（arguably, apparently, plausibly, relatively, typically）及谨慎动词（suggest, indicate, point to, appear to）控制断言边界。\par\vspace{4pt}
\textbf{Operational Rules}:
\begin{itemize}[leftmargin=1.5em, itemsep=1pt]
  \item 当存在未测混杂（unmeasured confounding）、小样本或仿真环境限制时，必须使用 may be attributable to / suggests that / supports a tentative conclusion。
  \item 禁止在经验研究中滥用 irrefutable, undeniable, inevitable, universal, conclusively prove；只有经核实的数学反例推翻全称命题时才可用 disprove / irrefutable。
\end{itemize}
\end{ConceptBox}
\begin{ExampleBox}{Authentic Examples \& Contrasts}
\textit{The observed discrepancy may be attributable to differences in measurement procedures.}\par\vspace{2pt}
\textit{A distributional forecast is arguably more informative when decisions depend on tail risk.}\par\vspace{2pt}
\end{ExampleBox}
\begin{TipBox}{Diagnostic Warning \& Common Errors}
\begin{itemize}[leftmargin=1.5em, itemsep=1pt]
  \item 在改写或总结时删去原句的 may / arguably / in this sample 等限定词，造成过度绝对化。
\end{itemize}
\end{TipBox}
\vspace{6pt}
\subsection{G015: 分词短语（现在分词与过去分词）作后置定语与状语及悬垂修饰防范}\label{sec:rule:G015}
\begin{ConceptBox}{Core Grammar Principle \& Theoretical Framework}
\textbf{Category}: \texttt{C7\_C11\_participles\_and\_dangling\_modifiers}\par\vspace{3pt}
\textbf{Principle}: 分词短语能将冗长的从句浓缩为紧凑的后置定语或状语；但句首分词短语的逻辑主语必须与主句主语完全一致，绝不可产生悬垂修饰（Dangling Modifier）。\par\vspace{4pt}
\textbf{Operational Rules}:
\begin{itemize}[leftmargin=1.5em, itemsep=1pt]
  \item 后置分词定语：records exceeding the limit (= that exceed...)；predictions compared with ground-truth labels (= that are compared...)。
  \item 句首分词状语主语一致铁律：Using a 48-step look-back window, the model reduces MSE（the model 可以 use window；但不能写 *Using a 48-step window, the MSE decreases，因为 MSE 不会自己 use window！）。
\end{itemize}
\end{ConceptBox}
\begin{ExampleBox}{Authentic Examples \& Contrasts}
\textit{Records exceeding the stated size limit were excluded.}\par\vspace{2pt}
\textit{Compared with earlier studies, this study used a larger sample.}\par\vspace{2pt}
\end{ExampleBox}
\begin{TipBox}{Diagnostic Warning \& Common Errors}
\begin{itemize}[leftmargin=1.5em, itemsep=1pt]
  \item 悬垂分词：*Training on synthetic data, the generalization error remained high（应改为 When the model was trained on synthetic data, ... 或 Training on synthetic data left the model with high generalization error）。
\end{itemize}
\end{TipBox}
\vspace{6pt}
\subsection{G016: 限制性/非限制性定语从句与“介词 + which/whom”精确修饰}\label{sec:rule:G016}
\begin{ConceptBox}{Core Grammar Principle \& Theoretical Framework}
\textbf{Category}: \texttt{C8\_relative\_clauses}\par\vspace{3pt}
\textbf{Principle}: 定语从句是构建精确后置定语、避免中式前置长定语的核心工具。必须严格区分限制性（无逗号，界定子集）与非限制性（有逗号，补充评注，不可用 that），并掌握介词前置结构。\par\vspace{4pt}
\textbf{Operational Rules}:
\begin{itemize}[leftmargin=1.5em, itemsep=1pt]
  \item 限制性定语从句缩窄先行词范围：The analogy breaks down where the two systems differ in the feature on which the comparison depends。
  \item 介词 + which 结构由从句动词或先行词搭配决定：the age at which, the point at which, the rate at which, the extent to which, a protocol by means of which, the assumptions on which the model rests。
\end{itemize}
\end{ConceptBox}
\begin{ExampleBox}{Authentic Examples \& Contrasts}
\textit{We devised a protocol by means of which records could be checked.}\par\vspace{2pt}
\textit{The study examines the degree to which access varies by region.}\par\vspace{2pt}
\end{ExampleBox}
\begin{TipBox}{Diagnostic Warning \& Common Errors}
\begin{itemize}[leftmargin=1.5em, itemsep=1pt]
  \item 非限制性定语从句逗号后误用 that。
  \item 介词 + which 遗漏固定介词：*the extent which the findings generalise（漏掉 to）。
\end{itemize}
\end{TipBox}
\vspace{6pt}
\subsection{G017: 名词性从句与同位语从句（the claim/assumption/fact that...）}\label{sec:rule:G017}
\begin{ConceptBox}{Core Grammar Principle \& Theoretical Framework}
\textbf{Category}: \texttt{C8\_noun\_clauses}\par\vspace{3pt}
\textbf{Principle}: 同位语从句（the claim/assumption/hypothesis/notion/contention/fact that + 完整句）用于精确展开抽象学术名词的内涵；介词后若需接完整陈述句，必须借助 the fact that 或 in that。\par\vspace{4pt}
\textbf{Operational Rules}:
\begin{itemize}[leftmargin=1.5em, itemsep=1pt]
  \item 同位语从句连接词 that 不作成分但不可省略，也不可换为 which（The contention that all data are interchangeable needs justification）。
  \item 介词（account for, despite, due to, stem from）不能直接接 that 从句，必须用 account for the fact that / despite the fact that / stem from the fact that。
\end{itemize}
\end{ConceptBox}
\begin{ExampleBox}{Authentic Examples \& Contrasts}
\textit{The model cannot account for the fact that rates differ across sites.}\par\vspace{2pt}
\textit{The scheme was adopted despite the fact that it was costly.}\par\vspace{2pt}
\end{ExampleBox}
\begin{TipBox}{Diagnostic Warning \& Common Errors}
\begin{itemize}[leftmargin=1.5em, itemsep=1pt]
  \item 介词直接接 that 从句：*despite that it was costly / *due to that several records were missing。
\end{itemize}
\end{TipBox}
\vspace{6pt}
\subsection{G018: 状语从句、让步/条件结构与介词短语的紧凑转换}\label{sec:rule:G018}
\begin{ConceptBox}{Core Grammar Principle \& Theoretical Framework}
\textbf{Category}: \texttt{C8\_C9\_adverbial\_clauses}\par\vspace{3pt}
\textbf{Principle}: 掌握从属连词（although, whereas, while, provided that, because, in that）、介词短语（despite, in spite of, due to, prior to, in light of, in the absence of）与连接副词（nevertheless, nonetheless, consequently, conversely）之间的句法转换。\par\vspace{4pt}
\textbf{Operational Rules}:
\begin{itemize}[leftmargin=1.5em, itemsep=1pt]
  \item 让步压缩：Although the approach has limitations $\to$ Albeit with limitations / Despite its limitations（严禁写 *despite of）。
  \item 条件限定：provided that / as long as + 完整分句 <$\to$ subject to / under + 名词短语。
  \item 对照结构：whereas / while 连接两个对称分句（The first model uses monthly data, whereas the second uses daily data）。
\end{itemize}
\end{ConceptBox}
\begin{ExampleBox}{Authentic Examples \& Contrasts}
\textit{The approach is informative, albeit with limitations.}\par\vspace{2pt}
\textit{The result applies provided that the assumptions hold.}\par\vspace{2pt}
\end{ExampleBox}
\begin{TipBox}{Diagnostic Warning \& Common Errors}
\begin{itemize}[leftmargin=1.5em, itemsep=1pt]
  \item 混淆介词与连词：*despite the sample is large。
  \item while 表对比时与时间含义发生歧义（在时序/并发场景中表对比优先用 whereas）。
\end{itemize}
\end{TipBox}
\vspace{6pt}
\section{Cluster 4: Comparisons, Parallelism \& Mechanics (G019--G024) / 比较论证、平行结构与标点规约}
\Needspace{6\baselineskip}
\subsection{G019: 比较结构、拉丁比较级（superior/prior to）与比较维度限定（in terms of）}\label{sec:rule:G019}
\begin{ConceptBox}{Core Grammar Principle \& Theoretical Framework}
\textbf{Category}: \texttt{C9\_comparison\_structures}\par\vspace{3pt}
\textbf{Principle}: 科学比较必须做到：(1) 比较主体语法同构，(2) 明确比较指标与条件（in terms of / with respect to / under comparable conditions），(3) 正确使用拉丁比较级与类比结构。\par\vspace{4pt}
\textbf{Operational Rules}:
\begin{itemize}[leftmargin=1.5em, itemsep=1pt]
  \item superior to, inferior to, prior to, analogous to, akin to 一律用介词 to，严禁用 than。
  \item 声明某方法更优（superior / optimal / outperform）时，必须限定评估指标与数据集范围，禁止无条件宣称“处处最优”。
\end{itemize}
\end{ConceptBox}
\begin{ExampleBox}{Authentic Examples \& Contrasts}
\textit{The method was superior to the baseline on this metric, but other outcomes were not assessed.}\par\vspace{2pt}
\textit{The procedure is akin to stratified sampling, although the allocation rule differs.}\par\vspace{2pt}
\end{ExampleBox}
\begin{TipBox}{Diagnostic Warning \& Common Errors}
\begin{itemize}[leftmargin=1.5em, itemsep=1pt]
  \item 误用 *superior than / *prior than。
  \item 省略比较基准（baseline）或评价维度。
\end{itemize}
\end{TipBox}
\vspace{6pt}
\subsection{G020: 语态的语篇决策：基于研究对象、旧信息前置与施事责任的主动/被动选择}\label{sec:rule:G020}
\begin{ConceptBox}{Core Grammar Principle \& Theoretical Framework}
\textbf{Category}: \texttt{D\_voice\_selection\_strategy}\par\vspace{3pt}
\textbf{Principle}: 破除“被动语态一定比主动语态更学术”的迷思。语态选择服务于：(1) 维持句间话题链连贯（把上句提到的受事对象作为本句被动主语，满足 Given-before-New），(2) 明确谁对设计决策或核验承担责任。\par\vspace{4pt}
\textbf{Operational Rules}:
\begin{itemize}[leftmargin=1.5em, itemsep=1pt]
  \item 当描述标准化数据清洗、传感器部署或客观流程，且受事对象是段落主题时，选用被动语态（The records are arranged in chronological order before the time-series analysis）。
  \item 当阐述本文核心创新设计、范围限定决策或强调多名标注者独立操作时，选用主动语态（We confine the analysis to...; Two analysts analysed the data independently; The algorithm exploits the sparsity of the matrix）。
\end{itemize}
\end{ConceptBox}
\begin{ExampleBox}{Authentic Examples \& Contrasts}
\textit{Passive for topic continuity: The predictions are compared with ground-truth labels that were checked by two annotators.}\par\vspace{2pt}
\textit{Active for clear methodological agency: We reformulate the constrained problem as an equivalent optimization problem.}\par\vspace{2pt}
\end{ExampleBox}
\begin{TipBox}{Diagnostic Warning \& Common Errors}
\begin{itemize}[leftmargin=1.5em, itemsep=1pt]
  \item 通篇机械使用被动语态导致语句冗长、施事责任模糊（如分不清是前人文献做的还是本文做的）。
\end{itemize}
\end{TipBox}
\vspace{6pt}
\subsection{G021: 句子改写核心规范：语义等价、限定守恒与消除四大改写偏差}\label{sec:rule:G021}
\begin{ConceptBox}{Core Grammar Principle \& Theoretical Framework}
\textbf{Category}: \texttt{F\_sentence\_rewriting\_rules}\par\vspace{3pt}
\textbf{Principle}: 句子改写必须在改变句法或词汇结构的同时，100\% 保持原句的事实范围、否定、时间、程度和认知确定性，并能准确诊断语法错误、表达不自然、语域不当与语义扭曲。\par\vspace{4pt}
\textbf{Operational Rules}:
\begin{itemize}[leftmargin=1.5em, itemsep=1pt]
  \item 限定守恒：原句的 provided that, in this sample, under Gaussian-kernel assumptions, may, arguably 等在改写后必须保留等价限定。
  \item 多路径自然改写：支持动词 <$\to$ 派生名词转换、从句 <$\to$ 介词短语转换、主动 <$\to$ 被动焦点转换，提供多个自然可接受版本。
\end{itemize}
\end{ConceptBox}
\begin{ExampleBox}{Authentic Examples \& Contrasts}
\textit{Source: Missing records account for the discrepancy between the totals.
Rewrite 1: The discrepancy between the totals is attributable to missing records.
Rewrite 2: Missing records explain why the two totals differ.}\par\vspace{2pt}
\end{ExampleBox}
\begin{TipBox}{Diagnostic Warning \& Common Errors}
\begin{itemize}[leftmargin=1.5em, itemsep=1pt]
  \item 在改写中偷换概念或将相关关系升级为因果结论。
\end{itemize}
\end{TipBox}
\vspace{6pt}
\subsection{G022: 并列、从属、平行结构（Parallelism）与对比对称性}\label{sec:rule:G022}
\begin{ConceptBox}{Core Grammar Principle \& Theoretical Framework}
\textbf{Category}: \texttt{C10\_parallelism\_and\_subordination}\par\vspace{3pt}
\textbf{Principle}: 由 and, or, rather than, as opposed to, as well as, both... and... 连接的成分必须在词性、句法层级与逻辑角色上严格平行。\par\vspace{4pt}
\textbf{Operational Rules}:
\begin{itemize}[leftmargin=1.5em, itemsep=1pt]
  \item 对比排除连接词（rather than, as opposed to）两侧结构必须平行：reporting observed values as opposed to predictions；measured an association rather than a causal effect。
  \item 分别对应结构（A and B ... X and Y, respectively）两侧元素个数与顺序必须严格一一对应。
\end{itemize}
\end{ConceptBox}
\begin{ExampleBox}{Authentic Examples \& Contrasts}
\textit{The table reports observed values as opposed to predictions.}\par\vspace{2pt}
\textit{The symbols x and y denote temperature and demand, respectively.}\par\vspace{2pt}
\end{ExampleBox}
\begin{TipBox}{Diagnostic Warning \& Common Errors}
\begin{itemize}[leftmargin=1.5em, itemsep=1pt]
  \item 两侧词性不对称：*combines quantitative measurements with interviewing participants（名词短语与动名词不对称）。
\end{itemize}
\end{TipBox}
\vspace{6pt}
\subsection{G023: 学术标点规范、分号/冒号逻辑功能与逗号粘连（Comma Splice）根治}\label{sec:rule:G023}
\begin{ConceptBox}{Core Grammar Principle \& Theoretical Framework}
\textbf{Category}: \texttt{C11\_punctuation\_and\_clause\_combining}\par\vspace{3pt}
\textbf{Principle}: 根治中国学生写作中最普遍的标点语法错误------逗号粘连（Comma Splice，即仅用逗号连接两个独立句子，或把 nevertheless, therefore, however 当连词用逗号连接两句）。\par\vspace{4pt}
\textbf{Operational Rules}:
\begin{itemize}[leftmargin=1.5em, itemsep=1pt]
  \item 连接副词（however, therefore, hence, nevertheless, nonetheless, consequently, moreover, furthermore, conversely, likewise）连接两个独立分句时，前句末尾必须用分号（;）或句号（.）！
  \item 前置复合修饰语加连字符（a small-scale study, an in-depth analysis, ground-truth labels, blue-collar occupations），后置作表语或状语时不加连字符（examine the issue in depth）。
\end{itemize}
\end{ConceptBox}
\begin{ExampleBox}{Authentic Examples \& Contrasts}
\textit{The observations are dependent; hence, the usual independence assumption does not apply.}\par\vspace{2pt}
\textit{The sample is small; nevertheless, the measurements identify a useful question for further study.}\par\vspace{2pt}
\end{ExampleBox}
\begin{TipBox}{Diagnostic Warning \& Common Errors}
\begin{itemize}[leftmargin=1.5em, itemsep=1pt]
  \item Comma Splice：*The sample is small, however, the protocol is reproducible（第一个逗号必须改为分号或句号）。
\end{itemize}
\end{TipBox}
\vspace{6pt}
\subsection{G024: 适度名词化（Nominalization）、复杂名词短语与动词活力平衡}\label{sec:rule:G024}
\begin{ConceptBox}{Core Grammar Principle \& Theoretical Framework}
\textbf{Category}: \texttt{C12\_nominalization\_and\_complex\_np}\par\vspace{3pt}
\textbf{Principle}: 名词化可将前句的过程浓缩为后句主语（如 The relaxation of this assumption...），实现紧凑衔接；但绝不能认为“名词化越多越学术”，过度名词化会导致句子僵化、头重脚轻。\par\vspace{4pt}
\textbf{Operational Rules}:
\begin{itemize}[leftmargin=1.5em, itemsep=1pt]
  \item 用“The + 派生名词 + of + 核心对象”打包已知过程作主语，谓语仍使用清晰有力的动词（The elimination of duplicate entries reduced the file size）。
  \item 当主语和宾语都已是长抽象名词时，中间谓语直接用强动词（affects, constrains, clarifies），不要写 *conducts an influence on。
\end{itemize}
\end{ConceptBox}
\begin{ExampleBox}{Authentic Examples \& Contrasts}
\textit{The emergence of new devices changed the monitoring procedure.}\par\vspace{2pt}
\textit{The quantification of uncertainty is central to this analysis.}\par\vspace{2pt}
\end{ExampleBox}
\begin{TipBox}{Diagnostic Warning \& Common Errors}
\begin{itemize}[leftmargin=1.5em, itemsep=1pt]
  \item 僵尸名词堆砌（Zombie Nouns）：*The implementation of the reduction of parameters was carried out（应直接写 We reduced the number of parameters）。
\end{itemize}
\end{TipBox}
\vspace{6pt}
\section{Cluster 5: Ambiguity, Causality \& Discourse Architecture (G025--G030) / 排歧机制、因果推断与篇章骨架}
\Needspace{6\baselineskip}
\subsection{G025: 句法歧义消除（介词短语附着、修辞辖域与比较基准排歧）}\label{sec:rule:G025}
\begin{ConceptBox}{Core Grammar Principle \& Theoretical Framework}
\textbf{Category}: \texttt{C11\_syntactic\_ambiguity\_resolution}\par\vspace{3pt}
\textbf{Principle}: 消除句子中的结构性多解（Structural Ambiguity），确保读者无需猜测修饰语究竟挂靠在哪个词上。\par\vspace{4pt}
\textbf{Operational Rules}:
\begin{itemize}[leftmargin=1.5em, itemsep=1pt]
  \item 调整介词短语位置或改写为定语从句，消除“V + NP + PP”中的 PP 附着歧义。
  \item 明确区分绝对差值（in absolute terms）与相对变化（in relative terms），并交代基准分母（denominator）。
\end{itemize}
\end{ConceptBox}
\begin{ExampleBox}{Authentic Examples \& Contrasts}
\textit{The absolute difference was reported alongside the relative change.}\par\vspace{2pt}
\textit{The revised definition is unambiguous.}\par\vspace{2pt}
\end{ExampleBox}
\begin{TipBox}{Diagnostic Warning \& Common Errors}
\begin{itemize}[leftmargin=1.5em, itemsep=1pt]
  \item 修饰语远离核心词导致语义滑稽或产生双重解释。
\end{itemize}
\end{TipBox}
\vspace{6pt}
\subsection{G026: 因果关系、统计相关、混杂（Confounding）与反事实（Counterfactual）表达边界}\label{sec:rule:G026}
\begin{ConceptBox}{Core Grammar Principle \& Theoretical Framework}
\textbf{Category}: \texttt{E\_G\_causality\_vs\_correlation}\par\vspace{3pt}
\textbf{Principle}: 在机器学习、流行病学、经济学与社会科学写作中，必须严格区分预测准确率/统计相关（prediction accuracy / correlation / association）与因果效应（causal effect / causality）。\par\vspace{4pt}
\textbf{Operational Rules}:
\begin{itemize}[leftmargin=1.5em, itemsep=1pt]
  \item 观察性关联用 correlate with, be associated with, coincide with；不能仅凭高预测精度或滞后项相关性（lagged correlation）声称确立了 causality。
  \item 讨论因果识别时，准确使用 confounding（混杂因素）、counterfactual outcome（反事实结局）、exogenous（外生的，需写明与扰动项关系）及 disentangle A from B（分离交织效应）。
\end{itemize}
\end{ConceptBox}
\begin{ExampleBox}{Authentic Examples \& Contrasts}
\textit{Predictive accuracy alone cannot establish causality.}\par\vspace{2pt}
\textit{Unmeasured confounding may bias an estimated treatment effect.}\par\vspace{2pt}
\end{ExampleBox}
\begin{TipBox}{Diagnostic Warning \& Common Errors}
\begin{itemize}[leftmargin=1.5em, itemsep=1pt]
  \item 在观察性数据或纯预测模型论文中滥用 cause, determine, prove a causal link。
\end{itemize}
\end{TipBox}
\vspace{6pt}
\subsection{G027: 文献综述、观点归属与报道动词（Reporting Verbs）的认知立场控制}\label{sec:rule:G027}
\begin{ConceptBox}{Core Grammar Principle \& Theoretical Framework}
\textbf{Category}: \texttt{E\_reporting\_verbs\_and\_attribution}\par\vspace{3pt}
\textbf{Principle}: 在 Related Work 与议论文引用中，通过精准选择报道动词（argue, claim, assert, contend, maintain, suggest, note, observe, point out, demonstrate, show, acknowledge, question, challenge）体现对前人工作的客观定位。\par\vspace{4pt}
\textbf{Operational Rules}:
\begin{itemize}[leftmargin=1.5em, itemsep=1pt]
  \item 区分事实呈现（show, demonstrate, reveal）、作者主观论断（claim, assert, contend, maintain）与审慎提议（suggest, propose, hypothesize）。
  \item According to the authors 仅用于引述他人或外部记录，绝不可写 *According to me。
\end{itemize}
\end{ConceptBox}
\begin{ExampleBox}{Authentic Examples \& Contrasts}
\textit{The authors argue that simpler baselines deserve attention.}\par\vspace{2pt}
\textit{The author asserts that the categories are exhaustive, whereas recent benchmarks challenge this assumption.}\par\vspace{2pt}
\end{ExampleBox}
\begin{TipBox}{Diagnostic Warning \& Common Errors}
\begin{itemize}[leftmargin=1.5em, itemsep=1pt]
  \item 用一成不变的 *Author A says that... 贯穿整个 Related Work 章节。
\end{itemize}
\end{TipBox}
\vspace{6pt}
\subsection{G028: 学术段落逻辑架构：主题句---机制解释---实证/实例支撑---范围限定}\label{sec:rule:G028}
\begin{ConceptBox}{Core Grammar Principle \& Theoretical Framework}
\textbf{Category}: \texttt{G\_paragraph\_structure\_and\_argumentation}\par\vspace{3pt}
\textbf{Principle}: 无论是顶会论文 Discussion 还是六级/雅思托福主体段，高质量段落都必须具备“清晰主题句（Topic Sentence）$\to$ 作用机制或理论解释（Mechanism/Warrant）$\to$ 具体数据或实例（Concrete Evidence/Instance）$\to$ 边界限定（Qualification）”四层逻辑，拒绝仅靠连接词填空。\par\vspace{4pt}
\textbf{Operational Rules}:
\begin{itemize}[leftmargin=1.5em, itemsep=1pt]
  \item 防范逻辑跳跃（Logical Leap）：从“某现象存在或某指标提升”到“应全面推广某政策/模型”之间，必须补充成本、可行性（feasibility）或分布偏移（distribution shift）的权衡（trade-off）。
  \item 正确区分对比连接词：on the contrary 仅用于否定刚提及的命题（“恰恰相反”），而对比两个不同群体的表现必须用 by contrast / in contrast / on the other hand。
\end{itemize}
\end{ConceptBox}
\begin{ExampleBox}{Authentic Examples \& Contrasts}
\textit{The report claims that every village was sampled. On the contrary, remote villages were excluded.}\par\vspace{2pt}
\textit{Urban use increased; by contrast, rural use remained stable.}\par\vspace{2pt}
\end{ExampleBox}
\begin{TipBox}{Diagnostic Warning \& Common Errors}
\begin{itemize}[leftmargin=1.5em, itemsep=1pt]
  \item 把 on the contrary 误当作 on the other hand / by contrast 使用。
  \item 段落只有空洞观点句和夸张形容词，缺乏具体机制解释与可核验实例。
\end{itemize}
\end{TipBox}
\vspace{6pt}
\subsection{G029: 语篇信息流推进（Given-New Contract）与跨句主题链构建}\label{sec:rule:G029}
\begin{ConceptBox}{Core Grammar Principle \& Theoretical Framework}
\textbf{Category}: \texttt{C12\_G\_information\_packaging}\par\vspace{3pt}
\textbf{Principle}: 地道英语段落之所以流畅自然、毫无生硬拼接感，核心在于遵守“旧信息在前（Given）、新信息在后（New）”的信息契约，用总结性指代或语态转换实现句间无缝咬合。\par\vspace{4pt}
\textbf{Operational Rules}:
\begin{itemize}[leftmargin=1.5em, itemsep=1pt]
  \item 链式推进（Linear Progression）：第 N 句句末引入的新概念（如 an anomaly in the sensor readings），在第 N+1 句句首打包为旧信息主语（This anomaly / Such discrepancies...）。
  \item 避免连续多句使用毫无关联的全新长主语，造成读者认知断层。
\end{itemize}
\end{ConceptBox}
\begin{ExampleBox}{Authentic Examples \& Contrasts}
\textit{An anomaly in the sensor readings prompted a check of the calibration procedure. This check revealed a drift in the reference voltage.}\par\vspace{2pt}
\end{ExampleBox}
\begin{TipBox}{Diagnostic Warning \& Common Errors}
\begin{itemize}[leftmargin=1.5em, itemsep=1pt]
  \item 把旧信息放在句末，而把长达 15 词的全新复杂概念硬塞在句首主语位置。
\end{itemize}
\end{TipBox}
\vspace{6pt}
\subsection{G030: 让步反驳（Concession-Rebuttal）、局限性陈述与外推边界（Extrapolation Limits）}\label{sec:rule:G030}
\begin{ConceptBox}{Core Grammar Principle \& Theoretical Framework}
\textbf{Category}: \texttt{E\_G\_concession\_rebuttal\_and\_limitations}\par\vspace{3pt}
\textbf{Principle}: 顶会论文 Limitations 章节、Rebuttal 及高分议论文的核心技能：坦诚承认方法或样本的约束条件（acknowledge constraints），同时清楚界定在什么范围内结论依然稳健（robust），并指出何时外推（extrapolation）会失效。\par\vspace{4pt}
\textbf{Operational Rules}:
\begin{itemize}[leftmargin=1.5em, itemsep=1pt]
  \item 区分内插（interpolation，在已观测范围内估计）与外推（extrapolation，向观测年龄、工况或分布范围外推断，需更强假设）。
  \item 用 While / Although / Despite 承认局限，用 nevertheless / aún / provided that 说明受控范围内的有效性与后续改进方向（warrant further investigation）。
\end{itemize}
\end{ConceptBox}
\begin{ExampleBox}{Authentic Examples \& Contrasts}
\textit{Extrapolation beyond the observed age range requires additional justification.}\par\vspace{2pt}
\textit{The discrepancy warrants further investigation, but it does not warrant a causal conclusion.}\par\vspace{2pt}
\end{ExampleBox}
\begin{TipBox}{Diagnostic Warning \& Common Errors}
\begin{itemize}[leftmargin=1.5em, itemsep=1pt]
  \item 要么完全隐瞒局限，要么自相矛盾地声称局限毫无影响。
\end{itemize}
\end{TipBox}
\vspace{6pt}
\section{Cluster 6: Anti-Defect Diagnostics \& Advanced Frontier Genres (G031--G036) / 病灶诊断与前沿文体适配}
\Needspace{6\baselineskip}
\subsection{G031: 句子成分精细剖析（主、谓、宾、表、定、状、补、同位语）与核心主干定位}\label{sec:rule:G031}
\begin{ConceptBox}{Core Grammar Principle \& Theoretical Framework}
\textbf{Category}: \texttt{C1\_sentence\_constituents\_analysis}\par\vspace{3pt}
\textbf{Principle}: 自然地道的英语写作建立在清晰的句子成分（Constituents）分工之上：主语（Subject，建议用简短核心名词短语）+ 谓语动词（Predicate Verb，承载核心动作或状态关系）+ 宾语/表语/补语（Object/Complement），而将限定范围、实验条件、时间窗口等次要信息合理分配给后置定语（Post-modifier）或状语（Adverbial）。\par\vspace{4pt}
\textbf{Operational Rules}:
\begin{itemize}[leftmargin=1.5em, itemsep=1pt]
  \item 主语（S）与谓语（V）的距离应尽可能紧凑：避免在主语核心词和谓语动词之间插入过长的非必要修饰成分，防止读者丢失句子主干。
  \item 区分定语（修饰名词，回答 which/what kind of）与状语（修饰动词或全句，回答 how/when/where/under what condition）：英語长定语一律后置（介词短语、分词短语、定语从句），切忌按中文习惯把长串修饰语全部堆在名词前面。
  \item 同位语（Appositive，如 a random walk, an archetypal example of a nonstationary process）用于紧随名词后补充等价定义或量纲说明，比单独另起一句更紧凑。
\end{itemize}
\end{ConceptBox}
\begin{ExampleBox}{Authentic Examples \& Contrasts}
\textit{[主语 The affine function f(x) = ax + b] [谓语 includes] [宾语 an intercept b].}\par\vspace{2pt}
\textit{[状语 When the relevant variances are finite and nonzero], [主语 lag-one Pearson autocorrelation] [谓语 measures] [宾语 the linear association] [后置定语 between observations one time step apart].}\par\vspace{2pt}
\end{ExampleBox}
\begin{TipBox}{Diagnostic Warning \& Common Errors}
\begin{itemize}[leftmargin=1.5em, itemsep=1pt]
  \item 中式前置长定语堆砌：*The collected from three rural clinics during the two-year pilot study data were noisy（应改为 The data collected from three rural clinics during the two-year pilot study were noisy）。
  \item 把状语介词短语误当成主语：*In Table 1 shows the ablation results（应为 Table 1 shows... 或 In Table 1, we report...）。
\end{itemize}
\end{TipBox}
\vspace{6pt}
\subsection{G032: 消除“头重脚轻（Top-Heavy Subject）”与遵循尾重原则（End-Weight Principle）}\label{sec:rule:G032}
\begin{ConceptBox}{Core Grammar Principle \& Theoretical Framework}
\textbf{Category}: \texttt{C10\_C12\_end\_weight\_and\_anti\_top\_heavy}\par\vspace{3pt}
\textbf{Principle}: 中国学习者与劣质机翻最常见的句法缺陷是“头重脚轻（Top-Heavy Sentence）”------主语长达十几个词，谓语却只是一个软弱的 is / occurs / exists。地道英语严格遵循“尾重原则（End-Weight）”：主语简短（旧信息），把复杂的从句、长列表或重信息放在谓语之后（句末焦点位置）。\par\vspace{4pt}
\textbf{Operational Rules}:
\begin{itemize}[leftmargin=1.5em, itemsep=1pt]
  \item 诊断头重脚轻：当主语超过 8--10 个词且包含从句或多层 of 短语，而谓语仅为 be/have/exist 时，必须重构句子。
  \item 重构方法 1（换用核心实体作主语+强动词）：把长主语中的行动者、模型或实验设为短主语，将原抽象名词还原为谓语动词（如把 *The investigation of how irregular sampling intervals affect autocorrelation estimation is important 改为 *We investigate how irregular sampling intervals affect autocorrelation estimation* 或 *Irregular sampling intervals can distort autocorrelation estimates*）。
  \item 重构方法 2（右移重成分或使用先行 it / 存在结构）：当需要评价某个复杂命题时，用 It is critical that... 或將状语条件移至句首/句末。
\end{itemize}
\end{ConceptBox}
\begin{ExampleBox}{Authentic Examples \& Contrasts}
\textit{头重脚轻病句：*Whether the learned policy can generalize from simulation to a physical robot with actuator delays and sensor noise is uncertain.*
地道改写：*It remains uncertain whether the learned policy can generalize from simulation to a physical robot with actuator delays and sensor noise.* 或 *Actuator delays and sensor noise make sim-to-real policy transfer uncertain.*}\par\vspace{2pt}
\textit{头重脚轻病句：*A discrepancy between the recorded timestamps and the sensor logs in the second rollout was found.*
地道改写：*The audit revealed a discrepancy between the recorded timestamps and the sensor logs in the second rollout.*}\par\vspace{2pt}
\end{ExampleBox}
\begin{TipBox}{Diagnostic Warning \& Common Errors}
\begin{itemize}[leftmargin=1.5em, itemsep=1pt]
  \item 主语塞满三四个 of 嵌套短语，导致读者读到句尾才看到一个单薄的 is。
  \item 用 Whether / That 名词性从句作长主语而谓语只有一个形容词。
\end{itemize}
\end{TipBox}
\vspace{6pt}
\subsection{G033: 去除“AI 味（Anti-AI-ese）”与过度修辞：用具体机制、指标与实例替代空洞夸张}\label{sec:rule:G033}
\begin{ConceptBox}{Core Grammar Principle \& Theoretical Framework}
\textbf{Category}: \texttt{E\_anti\_ai\_style\_and\_concreteness}\par\vspace{3pt}
\textbf{Principle}: 典型 AI 生成文本充斥着空洞的套话与浮夸副词/形容词（如 delve into, testament to, tapestry, underscores the paramount importance, revolutionize, unparalleled, groundbreaking），却缺乏具体技术细节。地道顶会论文（NeurIPS/ICLR/CoRL）与高分雅思/托福/六级写作追求“低修辞密度、高信息密度”：用具体动词、明确指标和真实实例说话。\par\vspace{4pt}
\textbf{Operational Rules}:
\begin{itemize}[leftmargin=1.5em, itemsep=1pt]
  \item 删除无实质信息的“清嗓子”套话（Throat-clearing）：少用 *It is worth noting that / It is undeniable that / In the realm of modern artificial intelligence*，直接陈述事实主干。
  \item 禁止“只有夸张形容词而不给实例或数据支撑”：不要写 *The model achieves an enormous, unparalleled improvement*，而要写 *The proposed parametrization reduces MSE by 8.4\% relative to the seasonal naive baseline on hourly electricity demand*。
  \item 优先选用朴素、精准的单音节或核心学术动词（examine, test, compare, constrain, decompose, bound, mitigate）替代故作高深的修辞词。
\end{itemize}
\end{ConceptBox}
\begin{ExampleBox}{Authentic Examples \& Contrasts}
\textit{AI 腔/空洞夸张：*This groundbreaking paradigm delves into the intricate tapestry of temporal dynamics and achieves unparalleled accuracy.*
地道学术改写：*This model decomposes the series into trend and seasonal components, achieving lower forecasting error than the linear baseline on three benchmarks.*}\par\vspace{2pt}
\textit{六级/雅思空洞句：*Technology has an enormous impact on education in many ways.*
地道具体展开：*Asynchronous learning platforms allow part-time students to review lecture recordings at their own pace, reducing scheduling conflicts.*}\par\vspace{2pt}
\end{ExampleBox}
\begin{TipBox}{Diagnostic Warning \& Common Errors}
\begin{itemize}[leftmargin=1.5em, itemsep=1pt]
  \item 堆砌 paramount, unprecedented, revolutionize, multifaceted 等宣传式词汇，违反客观学术语域。
  \item 只下宏大判断（如 *This method is very efficient*）而不给出具体的计算复杂度、内存占用或延迟数据。
\end{itemize}
\end{TipBox}
\vspace{6pt}
\subsection{G034: 四大常见表达病灶诊断：口语化、过度表达（Over-expression）、不精确表达与指代悬空}\label{sec:rule:G034}
\begin{ConceptBox}{Core Grammar Principle \& Theoretical Framework}
\textbf{Category}: \texttt{A\_E\_F\_four\_major\_writing\_defects}\par\vspace{3pt}
\textbf{Principle}: 针对中国学习者在论文、雅思托福及六级写作中的四大高频失分点建立系统修正机制：(1) 口语化表达混入书面语，(2) 过度表达（语义重复或情态过满），(3) 不精确表达（缺少比较基准、分母或时间单位），(4) 悬空指代。\par\vspace{4pt}
\textbf{Operational Rules}:
\begin{itemize}[leftmargin=1.5em, itemsep=1pt]
  \item 修正口语化（Colloquialism）：书面论文中将 a lot of 改为 substantial / a large number of；将 get 改为 obtain / acquire；将 big 改为 substantial / marked；将 do a study 改为 conduct / carry out a study；将 pretty up / dumb down 改为 refine the visual layout / oversimplify。
  \item 修正过度表达与语义冗余（Over-expression \& Pleonasm）：禁止写 *halve by 50\%*（halve 已含 50\%）、*mutually collaborate with each other*、*completely eliminate all*、*future forthcoming*。
  \item 修正不精确表达（Imprecision）：报告 lag（滞后）、rate（速率）或 relative change（相对变化）时，必须指明采样时间单位（sampling interval）、基线（baseline）与分母（denominator）。
\end{itemize}
\end{ConceptBox}
\begin{ExampleBox}{Authentic Examples \& Contrasts}
\textit{口语化病句：*We did a small study and got a lot of data, which pretty up our paper.*
修正：*We conducted a pilot study and collected a substantial amount of data to test the protocol.*}\par\vspace{2pt}
\textit{不精确病句：*The error dropped by 5\%.*
精确修正：*The classification error rate decreased by 5 percentage points (from 20\% to 15\%) relative to the baseline.*}\par\vspace{2pt}
\end{ExampleBox}
\begin{TipBox}{Diagnostic Warning \& Common Errors}
\begin{itemize}[leftmargin=1.5em, itemsep=1pt]
  \item 把口头演讲过渡语（如 *Getting back to the question / That is all very well, but*）直接写进正式期刊论文正文。
  \item 混淆相对百分比变化（per cent）与绝对百分点变化（percentage points）。
\end{itemize}
\end{TipBox}
\vspace{6pt}
\subsection{G035: 机器学习、时间序列、生成模型与具身智能（Embodied AI）的技术表达句法}\label{sec:rule:G035}
\begin{ConceptBox}{Core Grammar Principle \& Theoretical Framework}
\textbf{Category}: \texttt{E\_cs\_ai\_timeseries\_embodied\_syntax}\par\vspace{3pt}
\textbf{Principle}: 面向深度学习、时间序列预测、生成式建模与具身智能（Embodied AI）写作，训练如何用精准、克制、零 AI 腔的英语句式描述数学符号、分布特性、物理/几何约束、感知---动作闭环（perception-action loop）及实验消融。\par\vspace{4pt}
\textbf{Operational Rules}:
\begin{itemize}[leftmargin=1.5em, itemsep=1pt]
  \item 问题与符号定义：用 Let X denote... / Here, N denotes... / We parameterize the predictive distribution as... 保持全篇符号（notation）一致，严格区分标量（scalar）、向量（vector）与矩阵（matrix）。
  \item 时间序列与生成模型：严格区分平稳性（stationarity）与遍历性（ergodicity）、季节性（seasonality，固定日历周期）与周期波动（cyclicality，非固定时长）、概率校准（calibration，预测分布与观测频率一致）与锐度（sharpness，预测分布集中度）。
  \item 具身智能（Embodied AI）：准确使用 spatial orientation（空间方位）、trajectory（运动/策略轨迹）、velocity（速度矢量，含方向）、perturbation（外部扰动）、closed-loop feedback（闭环反馈）及 physical/kinematic constraints（物理与运动学约束），不把仿真结果（simulation）直接冒充真实物理部署（real-world deployment）。
\end{itemize}
\end{ConceptBox}
\begin{ExampleBox}{Authentic Examples \& Contrasts}
\textit{Time Series \& Generative: Forecast calibration concerns statistical consistency between predictive distributions and observations, whereas sharpness refers to the concentration of the predictive distribution.}\par\vspace{2pt}
\textit{Embodied AI: The controller enforces velocity and torque constraints while tracking the reference trajectory under external perturbations.}\par\vspace{2pt}
\end{ExampleBox}
\begin{TipBox}{Diagnostic Warning \& Common Errors}
\begin{itemize}[leftmargin=1.5em, itemsep=1pt]
  \item 把仿真环境（simulation）中的策略表现直接写成真实世界因果有效，忽略 sim-to-real discrepancy。
  \item 混淆 covariance（协方差，受量纲缩放影响）与 correlation（标准化相关系数），或混淆 isotonic（保序的）与 isotropic（各向同性的）。
\end{itemize}
\end{TipBox}
\vspace{6pt}
\subsection{G036: 跨场景文体适配：顶会论文、学术演讲（Tech Talk）、个人陈述（SoP）与六级/雅思托福高分议论文}\label{sec:rule:G036}
\begin{ConceptBox}{Core Grammar Principle \& Theoretical Framework}
\textbf{Category}: \texttt{E\_G\_multi\_genre\_adaptation}\par\vspace{3pt}
\textbf{Principle}: 同一核心词汇库在不同应用场景下需要灵活切换语域与句式结构：学术论文重实证边界与方法可复现；口头演讲（Presentation/Seminar）重听众引导与简洁路标；六级/雅思托福（CET-6/IELTS/TOEFL）重观点清晰、让步权衡与具体民生/科技/教育实例。\par\vspace{4pt}
\textbf{Operational Rules}:
\begin{itemize}[leftmargin=1.5em, itemsep=1pt]
  \item 顶会论文与审稿回复（Paper \& Rebuttal）：用 We thank the reviewer for pointing out this limitation; to address this concern, we conducted an additional sensitivity analysis... 礼貌、有据地回应质疑（take issue with / respond to / clarify）。
  \item 学术演讲与海报展示（Tech Talk \& Oral Presentation）：合理使用路标句（I would like to begin by explaining... / As this slide shows... / Let us now turn to... / I will leave five minutes for questions），句子宜短，避免听众无法跟上的多重嵌套从句。
  \item 六级与雅思托福写作（CET-6 \& IELTS/TOEFL Writing）：摒弃万能废话模板，采用“明确立场句（adopt a principled stance）+ 具体机制解释（account for / facilitate / mitigate）+ 落地实例（concrete example，如远程教育、可再生能源补贴、公共交通可及性）+ 让步权衡（albeit / on balance / outweigh）”。
\end{itemize}
\end{ConceptBox}
\begin{ExampleBox}{Authentic Examples \& Contrasts}
\textit{Tech Talk (Spoken Academic): As this slide shows, the naive baseline performs competitively when the look-back horizon is short; let us now turn to why complex models overfit noisy series.}\par\vspace{2pt}
\textit{IELTS/CET-6 Argumentative Essay: Subsidising public transport can facilitate access to education for rural students; however, municipal authorities must weigh this benefit against long-term infrastructure maintenance costs.}\par\vspace{2pt}
\end{ExampleBox}
\begin{TipBox}{Diagnostic Warning \& Common Errors}
\begin{itemize}[leftmargin=1.5em, itemsep=1pt]
  \item 在雅思/六级作文中使用生硬虚假的统计编造（如 *According to a 2025 Harvard survey, 93.7\% of people...*），而不是使用合理的机制论证与具体生活/政策实例。
  \item 在口头报告结束语与书面论文结论之间混用语域（如在论文结尾写 *That is all I have to say, thank you for listening*）。
\end{itemize}
\end{TipBox}
\vspace{6pt}
\chapter{Expressive Academic Sentence Templates (EP001--EP024) / 学术功能句型高阶模板库}\label{ch:templates24}
在学术论文的不同小节中，作者需要高效复用符合国际同行惯例的功能型表达句型（Expression Patterns）。本章系统呈现 24 组高分学术句型模板（EP001--EP024），涵盖从研究意图、量化对比、因果推断到前沿 AI 技术的全周期写作需求。

\section{Cluster A: Research Foundations, Scope \& Definitions (EP001--EP006) / 研究起步、范围定义与概念界定}
\Needspace{6\baselineskip}
\subsection{EP001: 明确陈述研究目标、问题定义与讨论边界}\label{sec:ep:EP001}
\begin{ConceptBox}{Template Formula \& Rhetorical Function}
\textbf{Function}: \texttt{stating\_aims\_and\_scope} \quad \textbf{Register}: \texttt{formal\_academic}\par\vspace{3pt}
\textbf{Template Formula}:\par\noindent\texttt{This study/paper examines [specific problem] and aims to [verb] [measurable objective] within [stated scope].}\par\vspace{4pt}
\textbf{Exposition}: 开门见山说明研究问题与范围，主语简短，直接用强动词 examines / addresses / investigates。
\end{ConceptBox}
\begin{ExampleBox}{Authentic Exemplars in Research Literature}
\textit{This paper addresses unequal access to educational resources.}\par\vspace{2pt}
\textit{The aim of this paper is to compare annotation rules.}\par\vspace{2pt}
\end{ExampleBox}
\begin{TipBox}{Academic Pitfall \& Reviewer Perspective}
\begin{itemize}[leftmargin=1.5em, itemsep=1pt]
  \item 使用冗长空洞的套话开头，或写 *address to a problem。
\end{itemize}
\end{TipBox}
\vspace{6pt}
\subsection{EP002: 精确区分可数样本量与不可数开销/证据}\label{sec:ep:EP002}
\begin{ConceptBox}{Template Formula \& Rhetorical Function}
\textbf{Function}: \texttt{describing\_methods\_and\_quantities} \quad \textbf{Register}: \texttt{general\_academic}\par\vspace{3pt}
\textbf{Template Formula}:\par\noindent\texttt{While a large number of [plural countable nouns] were collected, the procedure required only a small amount of [uncountable noun].}\par\vspace{4pt}
\textbf{Exposition}: 在数据集与实验开销描述中准确搭配可数与不可数量词。
\end{ConceptBox}
\begin{ExampleBox}{Authentic Exemplars in Research Literature}
\textit{A large number of records still require checking.}\par\vspace{2pt}
\textit{The protocol requires a small amount of additional information.}\par\vspace{2pt}
\end{ExampleBox}
\begin{TipBox}{Academic Pitfall \& Reviewer Perspective}
\begin{itemize}[leftmargin=1.5em, itemsep=1pt]
  \item *a great deal of participants / *an amount of records。
\end{itemize}
\end{TipBox}
\vspace{6pt}
\subsection{EP003: 客观归属文献观点或档案记录来源}\label{sec:ep:EP003}
\begin{ConceptBox}{Template Formula \& Rhetorical Function}
\textbf{Function}: \texttt{attributing\_sources\_and\_claims} \quad \textbf{Register}: \texttt{general\_academic}\par\vspace{3pt}
\textbf{Template Formula}:\par\noindent\texttt{According to [authors/records], [reported finding]; however, [methodological qualification].}\par\vspace{4pt}
\textbf{Exposition}: 准确区分外部文献说法与本文实证判断。
\end{ConceptBox}
\begin{ExampleBox}{Authentic Exemplars in Research Literature}
\textit{According to the authors, the evaluation uses synthetic sentences.}\par\vspace{2pt}
\textit{According to the records, the event began before sunset.}\par\vspace{2pt}
\end{ExampleBox}
\begin{TipBox}{Academic Pitfall \& Reviewer Perspective}
\begin{itemize}[leftmargin=1.5em, itemsep=1pt]
  \item *According to me（第一人称严禁使用）。
\end{itemize}
\end{TipBox}
\vspace{6pt}
\subsection{EP004: 界定概念的操作性定义与数学符号表示}\label{sec:ep:EP004}
\begin{ConceptBox}{Template Formula \& Rhetorical Function}
\textbf{Function}: \texttt{defining\_concepts\_and\_notation} \quad \textbf{Register}: \texttt{formal\_academic}\par\vspace{3pt}
\textbf{Template Formula}:\par\noindent\texttt{For the purposes of this analysis, we define [term] as [operational criterion], where [symbol] denotes [quantity].}\par\vspace{4pt}
\textbf{Exposition}: 用于将抽象概念转化为可观测、可复现的操作性指标或符号。
\end{ConceptBox}
\begin{ExampleBox}{Authentic Exemplars in Research Literature}
\textit{We define the outcome as the number of visits per week.}\par\vspace{2pt}
\textit{Here, N denotes the number of observations.}\par\vspace{2pt}
\end{ExampleBox}
\begin{TipBox}{Academic Pitfall \& Reviewer Perspective}
\begin{itemize}[leftmargin=1.5em, itemsep=1pt]
  \item 定义含糊或符号指代前后不一致。
\end{itemize}
\end{TipBox}
\vspace{6pt}
\subsection{EP005: 描述系统模块、数据集或分类体系的构成}\label{sec:ep:EP005}
\begin{ConceptBox}{Template Formula \& Rhetorical Function}
\textbf{Function}: \texttt{classification\_and\_composition} \quad \textbf{Register}: \texttt{formal\_academic}\par\vspace{3pt}
\textbf{Template Formula}:\par\noindent\texttt{[Whole System/Dataset] consists of / comprises [N] components, whereas [Subset] accounts for / constitutes [Proportion].}\par\vspace{4pt}
\textbf{Exposition}: 用主动语态的 consist of / comprise 表整体由部分组成，用 constitute / account for 表部分占整体的比例。
\end{ConceptBox}
\begin{ExampleBox}{Authentic Exemplars in Research Literature}
\textit{The questionnaire consists of three sections.}\par\vspace{2pt}
\textit{Missing responses account for a small proportion of the records.}\par\vspace{2pt}
\end{ExampleBox}
\begin{TipBox}{Academic Pitfall \& Reviewer Perspective}
\begin{itemize}[leftmargin=1.5em, itemsep=1pt]
  \item *is consisted of / *comprise of。
\end{itemize}
\end{TipBox}
\vspace{6pt}
\subsection{EP006: 描述实验执行、基线核查与可复现流程}\label{sec:ep:EP006}
\begin{ConceptBox}{Template Formula \& Rhetorical Function}
\textbf{Function}: \texttt{describing\_methods\_and\_procedures} \quad \textbf{Register}: \texttt{formal\_academic}\par\vspace{3pt}
\textbf{Template Formula}:\par\noindent\texttt{We conducted [experiment/survey] under controlled conditions and verified [readings/outputs] against [ground-truth/reference].}\par\vspace{4pt}
\textbf{Exposition}: 直接、清晰地交代实验操作与质量核验步骤，避免头重脚轻。
\end{ConceptBox}
\begin{ExampleBox}{Authentic Exemplars in Research Literature}
\textit{The team verified sensor readings against manual measurements.}\par\vspace{2pt}
\textit{Two analysts analysed the data independently.}\par\vspace{2pt}
\end{ExampleBox}
\begin{TipBox}{Academic Pitfall \& Reviewer Perspective}
\begin{itemize}[leftmargin=1.5em, itemsep=1pt]
  \item *make an experiment / *do a research。
\end{itemize}
\end{TipBox}
\vspace{6pt}
\section{Cluster B: Quantitative Analysis, Graphics \& Comparison (EP007--EP012) / 量化分析、图表汇报与多维对比}
\Needspace{6\baselineskip}
\subsection{EP007: 同时报告绝对差值与相对变化（区分 \% 与百分点）}\label{sec:ep:EP007}
\begin{ConceptBox}{Template Formula \& Rhetorical Function}
\textbf{Function}: \texttt{describing\_trends\_and\_changes} \quad \textbf{Register}: \texttt{formal\_academic}\par\vspace{3pt}
\textbf{Template Formula}:\par\noindent\texttt{[Metric] decreased/increased by [X] percentage points (from [A]\% to [B]\%), corresponding to a [Y]\% relative change.}\par\vspace{4pt}
\textbf{Exposition}: 消除百分比歧义，同时给出基线值、绝对差值与相对变化率。
\end{ConceptBox}
\begin{ExampleBox}{Authentic Exemplars in Research Literature}
\textit{The absolute difference was reported alongside the relative change.}\par\vspace{2pt}
\textit{A rise from 20\% to 25\% is five percentage points.}\par\vspace{2pt}
\end{ExampleBox}
\begin{TipBox}{Academic Pitfall \& Reviewer Perspective}
\begin{itemize}[leftmargin=1.5em, itemsep=1pt]
  \item 混淆 percentage points 与 per cent，或隐瞒基数分母。
\end{itemize}
\end{TipBox}
\vspace{6pt}
\subsection{EP008: 使用 be of + 抽象名词结构自然提升评价语域}\label{sec:ep:EP008}
\begin{ConceptBox}{Template Formula \& Rhetorical Function}
\textbf{Function}: \texttt{academic\_register\_and\_nominal\_style} \quad \textbf{Register}: \texttt{formal\_academic}\par\vspace{3pt}
\textbf{Template Formula}:\par\noindent\texttt{[Factor/Dataset] is of considerable/particular [importance/significance/value/relevance] to [target problem].}\par\vspace{4pt}
\textbf{Exposition}: 自然替换口语化的 be very important/useful。
\end{ConceptBox}
\begin{ExampleBox}{Authentic Exemplars in Research Literature}
\textit{The consent procedure is of particular significance to this study.}\par\vspace{2pt}
\textit{The evidence is of great relevance to the policy debate.}\par\vspace{2pt}
\end{ExampleBox}
\begin{TipBox}{Academic Pitfall \& Reviewer Perspective}
\begin{itemize}[leftmargin=1.5em, itemsep=1pt]
  \item *be of important（of 后必须接名词 importance）。
\end{itemize}
\end{TipBox}
\vspace{6pt}
\subsection{EP009: 谨慎归因与多因素解释（May Account For / Be Attributable To）}\label{sec:ep:EP009}
\begin{ConceptBox}{Template Formula \& Rhetorical Function}
\textbf{Function}: \texttt{hedging\_and\_attribution} \quad \textbf{Register}: \texttt{formal\_academic}\par\vspace{3pt}
\textbf{Template Formula}:\par\noindent\texttt{The observed discrepancy may be attributable to / stem from [Factor A] rather than [Factor B].}\par\vspace{4pt}
\textbf{Exposition}: 在分析实验误差或异常现象时，客观提出可能成因而不武断定论。
\end{ConceptBox}
\begin{ExampleBox}{Authentic Exemplars in Research Literature}
\textit{The observed discrepancy may be attributable to differences in measurement procedures.}\par\vspace{2pt}
\textit{Measurement error may account for part of the discrepancy.}\par\vspace{2pt}
\end{ExampleBox}
\begin{TipBox}{Academic Pitfall \& Reviewer Perspective}
\begin{itemize}[leftmargin=1.5em, itemsep=1pt]
  \item *attributable for / 把可能成因直接写成唯一确定原因。
\end{itemize}
\end{TipBox}
\vspace{6pt}
\subsection{EP010: 图表、坐标轴、斜率与分布形态的精准描述}\label{sec:ep:EP010}
\begin{ConceptBox}{Template Formula \& Rhetorical Function}
\textbf{Function}: \texttt{describing\_figures\_and\_tables} \quad \textbf{Register}: \texttt{formal\_academic}\par\vspace{3pt}
\textbf{Template Formula}:\par\noindent\texttt{In Figure [N], plotting [Y] against [X] reveals [pattern/slope], while the shaded band denotes [uncertainty interval].}\par\vspace{4pt}
\textbf{Exposition}: 用于论文结果部分与雅思小作文（Task 1），准确说明横纵轴、单位与误差带。
\end{ConceptBox}
\begin{ExampleBox}{Authentic Exemplars in Research Literature}
\textit{In the figure, a shaded band denotes the observation interval.}\par\vspace{2pt}
\textit{We plotted temperature against time.}\par\vspace{2pt}
\end{ExampleBox}
\begin{TipBox}{Academic Pitfall \& Reviewer Perspective}
\begin{itemize}[leftmargin=1.5em, itemsep=1pt]
  \item plot A against B 中颠倒纵轴 A 与横轴 B 的关系。
\end{itemize}
\end{TipBox}
\vspace{6pt}
\subsection{EP011: 多方法跨维度对比（In Terms Of / With Respect To / Whereas）}\label{sec:ep:EP011}
\begin{ConceptBox}{Template Formula \& Rhetorical Function}
\textbf{Function}: \texttt{comparing\_and\_contrasting} \quad \textbf{Register}: \texttt{formal\_academic}\par\vspace{3pt}
\textbf{Template Formula}:\par\noindent\texttt{The methods are compared in terms of [Metric A] and [Metric B]: [Method 1] excels in [Condition 1], whereas [Method 2] is more robust to [Condition 2].}\par\vspace{4pt}
\textbf{Exposition}: 有据、具体地对比不同模型或方案的优劣，拒绝一刀切吹嘘。
\end{ConceptBox}
\begin{ExampleBox}{Authentic Exemplars in Research Literature}
\textit{The methods are compared in terms of accuracy and computational cost.}\par\vspace{2pt}
\textit{The first model uses monthly data, whereas the second uses daily data.}\par\vspace{2pt}
\end{ExampleBox}
\begin{TipBox}{Academic Pitfall \& Reviewer Perspective}
\begin{itemize}[leftmargin=1.5em, itemsep=1pt]
  \item 比较主语悬垂或未指明比较维度。
\end{itemize}
\end{TipBox}
\vspace{6pt}
\subsection{EP012: 显式声明模型假设、研究范围与成立前提}\label{sec:ep:EP012}
\begin{ConceptBox}{Template Formula \& Rhetorical Function}
\textbf{Function}: \texttt{stating\_assumptions\_and\_constraints} \quad \textbf{Register}: \texttt{formal\_academic}\par\vspace{3pt}
\textbf{Template Formula}:\par\noindent\texttt{We confine the analysis to [domain/regime]; the bound holds provided that [regularity/stationarity assumption] is satisfied.}\par\vspace{4pt}
\textbf{Exposition}: 用于理论推导或实证边界说明，体现科学严谨性。
\end{ConceptBox}
\begin{ExampleBox}{Authentic Exemplars in Research Literature}
\textit{We confine the analysis to observations collected before the policy change.}\par\vspace{2pt}
\textit{A sufficient condition guarantees the stated result when that condition holds.}\par\vspace{2pt}
\end{ExampleBox}
\begin{TipBox}{Academic Pitfall \& Reviewer Perspective}
\begin{itemize}[leftmargin=1.5em, itemsep=1pt]
  \item 隐瞒模型成立的强假设。
\end{itemize}
\end{TipBox}
\vspace{6pt}
\section{Cluster C: Scientific Nuance, Causality \& Transitions (EP013--EP018) / 科学审慎、因果辨析与篇章过渡}
\Needspace{6\baselineskip}
\subsection{EP013: 严辨统计相关与因果机制（Correlation vs. Causality \& Confounding）}\label{sec:ep:EP013}
\begin{ConceptBox}{Template Formula \& Rhetorical Function}
\textbf{Function}: \texttt{distinguishing\_correlation\_and\_causation} \quad \textbf{Register}: \texttt{formal\_academic}\par\vspace{3pt}
\textbf{Template Formula}:\par\noindent\texttt{Although [Variable X] correlates with [Outcome Y], predictive accuracy alone cannot establish causality due to potential unmeasured confounding.}\par\vspace{4pt}
\textbf{Exposition}: 用于机器学习可解释性、医疗/社会科学数据分析，严格划清相关与因果的界限。
\end{ConceptBox}
\begin{ExampleBox}{Authentic Exemplars in Research Literature}
\textit{A correlation between two variables does not by itself establish a causal relationship.}\par\vspace{2pt}
\textit{Unmeasured confounding may bias an estimated treatment effect.}\par\vspace{2pt}
\end{ExampleBox}
\begin{TipBox}{Academic Pitfall \& Reviewer Perspective}
\begin{itemize}[leftmargin=1.5em, itemsep=1pt]
  \item 将特征重要性或滞后相关性直接称为因果作用。
\end{itemize}
\end{TipBox}
\vspace{6pt}
\subsection{EP014: 使用 respectively 与平行结构实现高密度信息配对}\label{sec:ep:EP014}
\begin{ConceptBox}{Template Formula \& Rhetorical Function}
\textbf{Function}: \texttt{precision\_and\_correspondence} \quad \textbf{Register}: \texttt{formal\_academic}\par\vspace{3pt}
\textbf{Template Formula}:\par\noindent\texttt{[Symbol/Group A] and [Symbol/Group B] denote/represent [Entity X] and [Entity Y], respectively.}\par\vspace{4pt}
\textbf{Exposition}: 在不增加从句负担的前提下清晰交代多组符号或实验参数的一一对应关系。
\end{ConceptBox}
\begin{ExampleBox}{Authentic Exemplars in Research Literature}
\textit{The symbols x and y denote temperature and demand, respectively.}\par\vspace{2pt}
\textit{The slope and intercept of f(x) = ax + b are a and b, respectively.}\par\vspace{2pt}
\end{ExampleBox}
\begin{TipBox}{Academic Pitfall \& Reviewer Perspective}
\begin{itemize}[leftmargin=1.5em, itemsep=1pt]
  \item 前后序列元素数量不匹配或顺序颠倒。
\end{itemize}
\end{TipBox}
\vspace{6pt}
\subsection{EP015: 章节过渡与Given-New指代衔接（This + 总结名词）}\label{sec:ep:EP015}
\begin{ConceptBox}{Template Formula \& Rhetorical Function}
\textbf{Function}: \texttt{discourse\_signposting} \quad \textbf{Register}: \texttt{formal\_academic}\par\vspace{3pt}
\textbf{Template Formula}:\par\noindent\texttt{This [discrepancy/limitation/pattern] motivates our next analysis, where we turn to [next topic] in greater detail.}\par\vspace{4pt}
\textbf{Exposition}: 用 This + 范畴名词总结上句结果并自然开启下一句或下一节。
\end{ConceptBox}
\begin{ExampleBox}{Authentic Exemplars in Research Literature}
\textit{An anomaly in the sensor readings prompted a check of the calibration procedure.}\par\vspace{2pt}
\textit{We now turn to the evaluation criteria.}\par\vspace{2pt}
\end{ExampleBox}
\begin{TipBox}{Academic Pitfall \& Reviewer Perspective}
\begin{itemize}[leftmargin=1.5em, itemsep=1pt]
  \item 使用指代不清的裸 This 作主语。
\end{itemize}
\end{TipBox}
\vspace{6pt}
\subsection{EP016: 多目标权衡（Trade-offs）与成本收益约束评估}\label{sec:ep:EP016}
\begin{ConceptBox}{Template Formula \& Rhetorical Function}
\textbf{Function}: \texttt{evaluating\_tradeoffs\_and\_arguments} \quad \textbf{Register}: \texttt{formal\_academic}\par\vspace{3pt}
\textbf{Template Formula}:\par\noindent\texttt{The design involves a trade-off between [computational cost / coverage] and [accuracy / latency]; on balance, the benefits outweigh the costs when [condition].}\par\vspace{4pt}
\textbf{Exposition}: 用于系统设计、模型复杂度分析以及六级/雅思托福利弊权衡类作文。
\end{ConceptBox}
\begin{ExampleBox}{Authentic Exemplars in Research Literature}
\textit{The proposal involves a trade-off between coverage and cost.}\par\vspace{2pt}
\textit{The expected benefits must outweigh the costs before the intervention is recommended.}\par\vspace{2pt}
\end{ExampleBox}
\begin{TipBox}{Academic Pitfall \& Reviewer Perspective}
\begin{itemize}[leftmargin=1.5em, itemsep=1pt]
  \item *outweigh than（outweigh 直接接宾语，不加 than）。
\end{itemize}
\end{TipBox}
\vspace{6pt}
\subsection{EP017: 坦诚说明研究局限与外推风险（Limitations \& Extrapolation）}\label{sec:ep:EP017}
\begin{ConceptBox}{Template Formula \& Rhetorical Function}
\textbf{Function}: \texttt{acknowledging\_limitations\_and\_extrapolation} \quad \textbf{Register}: \texttt{formal\_academic}\par\vspace{3pt}
\textbf{Template Formula}:\par\noindent\texttt{We acknowledge that [dataset/simulation limitation]; therefore, one should not extrapolate these findings to [unseen regime] without further validation.}\par\vspace{4pt}
\textbf{Exposition}: 用于论文 Limitations 小节，体现成熟严谨的科学态度。
\end{ConceptBox}
\begin{ExampleBox}{Authentic Exemplars in Research Literature}
\textit{The authors acknowledge that the sample excludes rural clinics.}\par\vspace{2pt}
\textit{We should not extrapolate from a single district to the entire country without justification.}\par\vspace{2pt}
\end{ExampleBox}
\begin{TipBox}{Academic Pitfall \& Reviewer Perspective}
\begin{itemize}[leftmargin=1.5em, itemsep=1pt]
  \item 用空洞借口掩盖核心样本偏差。
\end{itemize}
\end{TipBox}
\vspace{6pt}
\subsection{EP018: 全篇总结与后续研究动议（Conclusion \& Future Directions）}\label{sec:ep:EP018}
\begin{ConceptBox}{Template Formula \& Rhetorical Function}
\textbf{Function}: \texttt{concluding\_and\_motivating\_research} \quad \textbf{Register}: \texttt{formal\_academic}\par\vspace{3pt}
\textbf{Template Formula}:\par\noindent\texttt{In summary, the empirical evidence supports [claim] under [tested conditions], while the remaining anomalies warrant further investigation in [future setting].}\par\vspace{4pt}
\textbf{Exposition}: 收束全文核心贡献，不夸大、不遗漏边界条件。
\end{ConceptBox}
\begin{ExampleBox}{Authentic Exemplars in Research Literature}
\textit{The discrepancy warrants further investigation, but it does not warrant a causal conclusion.}\par\vspace{2pt}
\textit{The authors conclude that further external validation is needed.}\par\vspace{2pt}
\end{ExampleBox}
\begin{TipBox}{Academic Pitfall \& Reviewer Perspective}
\begin{itemize}[leftmargin=1.5em, itemsep=1pt]
  \item *warrant to do（应为 warrant + 名词/doing）。
\end{itemize}
\end{TipBox}
\vspace{6pt}
\section{Cluster D: Frontier Domains, Modern AI \& High-Stakes Genres (EP019--EP024) / 前沿技术、人工智能与高阶文体}
\Needspace{6\baselineskip}
\subsection{EP019: 机器学习与时间序列：问题定义、滞后项与符号约定}\label{sec:ep:EP019}
\begin{ConceptBox}{Template Formula \& Rhetorical Function}
\textbf{Function}: \texttt{ml\_timeseries\_problem\_formulation} \quad \textbf{Register}: \texttt{technical\_specialized}\par\vspace{3pt}
\textbf{Template Formula}:\par\noindent\texttt{Let [symbol] denote the observation at time t; we decompose the series into [components] and evaluate predictions over horizon H.}\par\vspace{4pt}
\textbf{Exposition}: 用于机器学习与时间序列论文的 Problem Formulation 部分，清晰界定输入历史窗口、滞后项（lagged observations）与预测步长。
\end{ConceptBox}
\begin{ExampleBox}{Authentic Exemplars in Research Literature}
\textit{Here, N denotes the number of observations, and a lag of one step links the current state to the preceding observation.}\par\vspace{2pt}
\textit{Time-series decomposition separates a series into components such as trend, seasonality and remainder.}\par\vspace{2pt}
\end{ExampleBox}
\begin{TipBox}{Academic Pitfall \& Reviewer Perspective}
\begin{itemize}[leftmargin=1.5em, itemsep=1pt]
  \item 未说明采样间隔（sampling interval）就直接声称某个 lag 代表固定物理时长。
\end{itemize}
\end{TipBox}
\vspace{6pt}
\subsection{EP020: 生成模型与概率预测：分布校准、不确定性量化与尾部风险}\label{sec:ep:EP020}
\begin{ConceptBox}{Template Formula \& Rhetorical Function}
\textbf{Function}: \texttt{generative\_and\_probabilistic\_modeling} \quad \textbf{Register}: \texttt{technical\_specialized}\par\vspace{3pt}
\textbf{Template Formula}:\par\noindent\texttt{A probabilistic/generative model assigns probabilities to possible outcomes, allowing us to quantify uncertainty and evaluate both calibration and sharpness.}\par\vspace{4pt}
\textbf{Exposition}: 用于生成模型（Generative Models）与概率预测论文，阐述为何需要学习完整条件分布而非单一点估计（point estimate）。
\end{ConceptBox}
\begin{ExampleBox}{Authentic Exemplars in Research Literature}
\textit{A distributional forecast is arguably more informative when decisions depend on tail risk.}\par\vspace{2pt}
\textit{Forecast calibration concerns statistical consistency between forecasts and observations.}\par\vspace{2pt}
\end{ExampleBox}
\begin{TipBox}{Academic Pitfall \& Reviewer Perspective}
\begin{itemize}[leftmargin=1.5em, itemsep=1pt]
  \item 将高集中度（sharpness）误认为已经过概率校准（calibrated）。
\end{itemize}
\end{TipBox}
\vspace{6pt}
\subsection{EP021: 具身智能（Embodied AI）与机器人：空间表征、轨迹跟踪与物理约束}\label{sec:ep:EP021}
\begin{ConceptBox}{Template Formula \& Rhetorical Function}
\textbf{Function}: \texttt{embodied\_ai\_and\_robotics} \quad \textbf{Register}: \texttt{technical\_specialized}\par\vspace{3pt}
\textbf{Template Formula}:\par\noindent\texttt{To bridge the discrepancy between simulation and real-world deployment, the policy enforces [physical/monotonicity constraints] while adapting to [sensor noise / external perturbations].}\par\vspace{4pt}
\textbf{Exposition}: 用于具身智能、机器人控制与视觉导航写作，准确表达仿真到真实迁移（sim-to-real）、空间方位（spatial orientation）与速度/运动学约束。
\end{ConceptBox}
\begin{ExampleBox}{Authentic Exemplars in Research Literature}
\textit{The optimization procedure enforces a non-negativity constraint on the weights while tracking the trajectory.}\par\vspace{2pt}
\textit{The sensitivity analysis examines the effect of a small perturbation to the input sensor readings.}\par\vspace{2pt}
\end{ExampleBox}
\begin{TipBox}{Academic Pitfall \& Reviewer Perspective}
\begin{itemize}[leftmargin=1.5em, itemsep=1pt]
  \item 把仿真器（simulation）中的理想合成观测（synthetic data）直接等同于真实物理传感器的含噪观测。
\end{itemize}
\end{TipBox}
\vspace{6pt}
\subsection{EP022: 消融实验（Ablation Study）与基线公平对比（Fair Baseline Comparison）}\label{sec:ep:EP022}
\begin{ConceptBox}{Template Formula \& Rhetorical Function}
\textbf{Function}: \texttt{ablation\_and\_experimental\_evaluation} \quad \textbf{Register}: \texttt{formal\_academic}\par\vspace{3pt}
\textbf{Template Formula}:\par\noindent\texttt{To disentangle the contribution of [Component A] from [Component B], we compare the full model against a conventional baseline under comparable conditions.}\par\vspace{4pt}
\textbf{Exposition}: 用于实验与消融分析章节，严谨剥离各模块的独立贡献，避免只挑有利指标汇报（selective reporting）。
\end{ConceptBox}
\begin{ExampleBox}{Authentic Exemplars in Research Literature}
\textit{Every proposed method should be compared against a relevant baseline under comparable conditions.}\par\vspace{2pt}
\textit{The study design aims to disentangle the effect of training from prior experience.}\par\vspace{2pt}
\end{ExampleBox}
\begin{TipBox}{Academic Pitfall \& Reviewer Perspective}
\begin{itemize}[leftmargin=1.5em, itemsep=1pt]
  \item 使用 unparallelled / groundbreaking 等空洞吹嘘词代替具体的消融实验差值与误差棒说明。
\end{itemize}
\end{TipBox}
\vspace{6pt}
\subsection{EP023: 六级与雅思托福（CET-6 / IELTS / TOEFL）议论文：去模板化“主张---机制---具体实例---让步”结构}\label{sec:ep:EP023}
\begin{ConceptBox}{Template Formula \& Rhetorical Function}
\textbf{Function}: \texttt{cet6\_ielts\_toefl\_essay\_writing} \quad \textbf{Register}: \texttt{general\_academic}\par\vspace{3pt}
\textbf{Template Formula}:\par\noindent\texttt{While [Concession: competing view/cost], [Main Claim: policy/technology] can [Verb: facilitate/mitigate] [Outcome] by [Concrete Mechanism], as illustrated by [Specific Real-World Case].}\par\vspace{4pt}
\textbf{Exposition}: 专为六级、雅思 Task 2 与托福独立写作设计：告别陈腐套话，用一句话完成让步对比，紧接着给出具体作用机制与可信实例。
\end{ConceptBox}
\begin{ExampleBox}{Authentic Exemplars in Research Literature}
\textit{While expanding public transport requires substantial infrastructure investment, subsidising bus fares can facilitate access to education for low-income students by reducing daily commuting costs.}\par\vspace{2pt}
\textit{Asynchronous online courses allow working adults to study at their own pace; nevertheless, they cannot wholly replace face-to-face seminars that stimulate spontaneous debate.}\par\vspace{2pt}
\end{ExampleBox}
\begin{TipBox}{Academic Pitfall \& Reviewer Perspective}
\begin{itemize}[leftmargin=1.5em, itemsep=1pt]
  \item 只写夸张口号（如 *Online education is a double-edged sword and plays a vital role*）而不展开具体原因机制与实例。
\end{itemize}
\end{TipBox}
\vspace{6pt}
\subsection{EP024: 学术演讲（Tech Talk）、研究陈述（SoP）与审稿回复（Rebuttal）的自然衔接}\label{sec:ep:EP024}
\begin{ConceptBox}{Template Formula \& Rhetorical Function}
\textbf{Function}: \texttt{oral\_presentation\_sop\_and\_rebuttal} \quad \textbf{Register}: \texttt{spoken\_academic}\par\vspace{3pt}
\textbf{Template Formula}:\par\noindent\texttt{In response to the reviewer's question regarding [X], we have clarified [Y] and added an experiment on [Z]; as this slide/table shows, the conclusion remains robust.}\par\vspace{4pt}
\textbf{Exposition}: 用于国际会议口头报告、博士/教职申请研究陈述（Research Statement）及同行评审回复，做到礼貌、直接、有理有据。
\end{ConceptBox}
\begin{ExampleBox}{Authentic Exemplars in Research Literature}
\textit{I would like to begin by explaining the sampling procedure, and I will allow time for questions at the end.}\par\vspace{2pt}
\textit{The revision addresses the limitation explicitly and incorporates feedback from the pilot evaluation.}\par\vspace{2pt}
\end{ExampleBox}
\begin{TipBox}{Academic Pitfall \& Reviewer Perspective}
\begin{itemize}[leftmargin=1.5em, itemsep=1pt]
  \item 在 Rebuttal 中情绪化反驳（hostile remarks）或回避审稿人指出的核心假设缺陷。
\end{itemize}
\end{TipBox}
\vspace{6pt}